\documentclass[11pt]{article}
\usepackage[margin=24mm]{geometry}
\usepackage{amsmath,amssymb,amsthm,mathtools,booktabs,hyperref}
\usepackage{graphicx,microtype}
\setlength{\emergencystretch}{2em}
\hypersetup{colorlinks=true,linkcolor=blue,urlcolor=blue,citecolor=blue}
\newtheorem{theorem}{Theorem}
\newtheorem{corollary}[theorem]{Corollary}
\newtheorem{proposition}[theorem]{Proposition}
\newcommand{\Herm}{\operatorname{Herm}}
\newcommand{\diag}{\operatorname{diag}}
\newcommand{\one}{\mathbf{1}}
\newcommand{\ii}{\mathrm{i}}
\newcommand{\status}[1]{\par\smallskip\noindent\textbf{Status: \texttt{#1}.}\ }

\title{Beyond Nodal Damping:\\Collective Interaction Bounds for Delay-Aware SG--GFL Co-Design}
\author{Jorge Luis Mayorga Taborda\\Universidad de los Andes, Colombia}
\date{3 October 2026}
\begin{document}
\maketitle
\begin{abstract}
Local PLL tuning and measurement delay act through the whole electrical
network. Starting from the exact delayed linearization, we show that
changing one PLL delay produces a rank-one update of the physical
torque--speed impedance. With independent propagation paths and a
nonzero update, its Hermitian damping increment has one positive
and one negative eigenvalue: latency redistributes harmonic damping
across collective directions.
For a fixed-equilibrium SG/GFL sharing contract, retaining network
algebraic variables also gives an affine descriptor perturbation.
A contour identity yields finite spectral-cluster changes, their full
parameter Hessian, and an explicit closed-walk remainder.
Necessary cluster conditions lead to linear and semidefinite outer
relaxations of secure replacement, conditional on verified contour
and remainder bounds.
In IEEE--39, all 60 declared harmonic tests confirm the finite-delay
identity numerically. Three frozen algebraic points give walk remainder
bounds 18--46 times smaller than a scalar-norm bound. Complex ball
arithmetic certifies one root cluster throughout a declared parameter box.
A preserved analytical gain construction at 40 ms delay yields 87.6\%
GFL and passes five nonlinear events. A new matched Julia comparison
finds that unchanged gains also pass all five, with small opposing
frequency and RoCoF changes. The results establish conditional
interaction laws and a feasible construction; they do not establish
replacement superiority or a physically limiting maximum.
\end{abstract}
\noindent\textbf{Index terms:} grid-following converters, PLL delay,
descriptor systems, nonlinear eigenvalues, collective interactions,
regional replacement bounds.

\section{Introduction and research question}
Replacing synchronous generation changes inertia, electrical coupling
and converter synchronization simultaneously. Retuning each PLL against
one operating point can preserve an assigned oscillatory mode while
moving other modes toward a security boundary. The design question is
therefore quantitative: within a declared replacement and gain region,
how much replacement can any admissible local PLL tuning support,
at fixed physical delays and under a stated security contract?

Detailed low-inertia models already connect converter control settings
to collective modal behavior \cite{markovic}. Grid structure and PLL
synchronization are also established concerns \cite{huang}.
Stability-constrained planning incorporates frequency and grid-strength
constraints \cite{chuteng}; recent expansion-planning work explicitly
addresses converter-driven modes \cite{jia}. Thus neither PLL tuning,
delay-induced instability, nor the use of network information is by
itself a new contribution.

The methodological distinction examined here is a \emph{necessary
finite regional bound across all allowed local gains}, rather than
only a successful tuning or a fixed-pattern assignment obstruction.
Descriptor lifting \cite{peaucelle}, nonlinear-eigenvalue contour
localization \cite{bindel}, contour moments \cite{beyn}, structured
DDE robustness \cite{borgioli}, and fixed-structure delayed-controller
optimization \cite{michielsdesign} are antecedents, not discoveries
of this work. The proposed synthesis specializes them to the actual
SG/GFL sharing model and makes inter-site curvature and remainder
terms explicit.

The manuscript establishes four conditional mathematical results:
a finite physical damping-redistribution identity, an affine full-model
intervention representation, a cluster curvature/closed-walk bound,
and a controller-bounded necessary replacement relaxation. Its numerical evidence includes identity and enclosure tests of the
bound machinery, a frozen full-model feasible retuning step, and a
matched five-event fixed-gain follow-up. An informative
IEEE--39 replacement ceiling below the parameter-box ceiling has not
been demonstrated. This limitation prevents a claim of superior
replacement capacity or certified near-optimality.

\paragraph{Organization.}
The main text states the physical contract, derives the finite
regional bounds, and reports their evidence and limitations.
Appendices preserve the physical damping derivation, frequency-moment
law, nonlinear RoCoF argument, exact gain map, and local design
constructions. These auxiliary results have their own assumptions;
they are not combined into an unstated nonlinear stability theorem.
\section{Physical model and optimization contract}\label{sec:model}
With fixed device support, write the nonlinear retarded DAE as
\begin{align}
 E(x,\rho)\dot x&=f_0(x,y;\rho)+\sum_i b_i(K_{p,i},K_{i,i})e_i(t-\tau_i),\\
 0&=g(x,y;\rho,d),\qquad
 e_i=-\sin\theta_i\,v_{r,i}+\cos\theta_i\,v_{i,i}. \label{eq:dae}
\end{align}
The exported reduced state convention has $E=I$. Physical inertia is
recovered by consistent rescaling of SG swing rows, not by mixing row
conventions. The implemented PLL is exactly
\begin{equation}
 \dot\theta_i=\omega_{p,i},\quad
 t_{f,i}\dot\omega_{p,i}=\xi_i+K_{p,i}e_i(t-\tau_i)-\omega_{p,i},
 \quad \dot\xi_i=K_{i,i}e_i(t-\tau_i). \label{eq:pll}
\end{equation}
Here $t_{f,i}=1/(300\,2\pi)$ s is the existing filter constant, distinct
from pure delay $\tau_i$. Other converter equations use $\theta_i$ but
do not directly use $\omega_{p,i}$ or $\xi_i$ in this implementation.
No current-loop, setpoint or DC signal is delayed by this contract.

Use initialized generator dispatch and rating:
\[
 \epsilon_i=1-\rho_i,\quad P_{SG,i}=\epsilon_iP_i^0,\quad
 P_{GFL,i}=\rho_iP_i^0,\quad M_i=2H_iS_i^0\epsilon_i,\quad
 J_{SG}=\sum_iP_i^0\epsilon_i .
\]
$S_i^0$ and $P_i^0$ are different quantities. Replacement changes
electrical injection/coupling as well as inertia. State removal at
$\rho_i=0$ or $1$ requires another support chart.

For consistent per-unit phasors $P_i+\ii Q_i=V_iI_i^*$.
The lossy network and voltage/current/DC dynamics remain inside $g,f$.
If $g_y$ is nonsingular, algebraic elimination gives
\[
 \delta y=-g_y^{-1}(g_x\delta x+g_d\delta d),\qquad
 A_{\rm inst}=f_x-f_yg_y^{-1}g_x
\]
in ODE coordinates, with the analogous substitution in each detector.
Thus P/Q and voltage effects enter the actual characteristic operator;
a lossless Laplacian does not replace the AC Jacobian.
No converter current limiter is implemented, so no limiter safety is claimed.

Sources inspected include:
\begin{itemize}
\item \path{experiments/nonlinear_codesign_20261001/ReducedDAE.jl};
\item \path{src/bnd_model_expN/PDExactDesignN.jl};
\item \path{experiments/graph_gsp_codesign_20261003/DelayedEvents.jl}.
\end{itemize}
The intended DC convention is physical supply in \path{ReducedDAE.jl};
the older helper by itself uses a different convention.


For a compact support-preserving region $\mathcal D$, define
\[
 P^\star_{\mathcal D}(\tau)=
 \sup_{(\rho,K_p,K_i)\in\mathcal D}
 (P^0)^T\rho
\]
over designs satisfying the complete required nonrotational spectral
margin and every declared nonlinear-event constraint.
The delays $\tau_i$ are exogenous. A finite event set is a specified
security test, not robustness to arbitrary disturbances.
The initialized dispatch $(P^0)^T\mathbf1$ is the denominator for
the replacement fraction; net load and installed capacity are not
substituted for it.

All identities involving fixed exported matrices apply to the
repository's fixed-support, fixed-equilibrium current-sharing chart.
The internal per-unit device equations and trims are held fixed
while their terminal-current contributions are scaled. Physical
rating/inertia interpretation must use that same normalization.
Redispatch, topology changes, removal of device states at endpoints,
or different nonlinear converter models require a new derivation
or explicitly bounded model terms.
\section{A finite regional replacement bound with moving modal patterns}
\label{sec:regional}

A regional necessary bound must permit the modal patterns to move
with the design. The following construction uses complete analytic
characteristic matrices and bounds root clusters on fixed contours.
It does not impose the eigenvector compatibility conditions used
by the separate gain-construction method.

\subsection{A finite identity that survives internal root collisions}
\status{EXACT\_IDENTITY}
Use the full finite-dimensional characteristic matrix
$\Delta(s,p)$ of the linearized DDE, with fixed delays.
It is holomorphic in $s$ on and inside each positively oriented,
piecewise smooth simple contour $\Gamma$ used below.
Do not use a meromorphic Schur complement without accounting for its poles.
Exclude the rotational root from these contours. In the full state chart
this avoids requiring a parameter-independent gauge-deflation matrix.

Let $p=p_0+d$ lie in a declared compact design region $\mathcal D$ and
define $\mathcal D_d=\{d:p_0+d\in\mathcal D\}$, and let $\Delta_0(s)=\Delta(s,p_0)$ be nonsingular on $\Gamma$.
Suppose a factorization
\[
 \Delta(s,p)-\Delta_0(s)=\mathcal U(s)\mathcal V(s,d),\qquad
 H(s,d)=\mathcal V(s,d)\Delta_0(s)^{-1}\mathcal U(s)
\]
is available with $H\in\mathbb C^{r\times r}$, $\mathcal V(s,0)=0$.
All dimensions and coordinates remain fixed on $\mathcal D$.
Assume the \emph{uniform} bound
\begin{equation}
 \sup_{d\in\mathcal D_d,\,s\in\Gamma}\|H(s,d)\|_2
 \le\kappa_\Gamma<1 . \label{eq:regional-kappa}
\end{equation}
The determinant lemma and invertibility of $I+tH$, $0\le t\le1$,
preserve the number $m_\Gamma>0$ of enclosed roots, counting multiplicity.
Define their mean $z_\Gamma(p)=m_\Gamma^{-1}
\sum_{\lambda\ {\rm inside}\ \Gamma}\lambda$.
Then
\begin{equation}
 \boxed{z_\Gamma(p)-z_\Gamma(p_0)
 =-\frac{1}{2\pi\ii m_\Gamma}
       \oint_\Gamma \operatorname{tr}\log(I+H(s,d))\,ds .}
 \label{eq:regional-log}
\end{equation}
\begin{proof}
The determinant ratio on the contour is $\det(I+H)$.
The norm-convergent series
$\operatorname{tr}\log(I+H)=
\sum_{\ell\ge1}(-1)^{\ell+1}\operatorname{tr}(H^\ell)/\ell$
chooses a continuous, single-valued logarithm along the closed contour.
The argument principle gives the difference of root sums as
$(2\pi\ii)^{-1}\oint s\,d\log\det(I+H)$.
Integration by parts gives \eqref{eq:regional-log}.
This only uses regularity on the contour and constant enclosed count;
roots may collide or become defective inside it.
\end{proof}

If every required nonrotational physical root satisfies $\Re\lambda\le-\sigma$, then necessarily
\[
 g_\Gamma(p):=\Re z_\Gamma(p)+\sigma\le0.
\]
The converse is false: a stable mean can hide an unstable member.
Several disjoint contours can sharpen the necessary conditions.
One-root contours recover individual sensitivities where regular.
A finite collection is sufficient for a \emph{necessary relaxation},
not for a full-spectrum feasibility claim.

\subsection{An explicit uniform remainder}
\status{PROVED\_REDUCED\_MODEL}
Write $H=H_1+H_{\rm rem}$ with $H_1$ linear in $d$, and suppose
\[
 \sup_{\Gamma,\mathcal D_d}\|H_{\rm rem}\|_*
 \le e_\Gamma,\qquad \|\cdot\|_* \text{ is the nuclear norm}.
\]
Let $\ell_\Gamma$ be the contour length and put
$\varphi(x)=-\log(1-x)-x$, $0\le x<1$.
The affine prediction and its remainder satisfy
\begin{align}
 g_\Gamma(p)&=g_\Gamma(p_0)+J_\Gamma d+\mathcal R_\Gamma(d),\\
 J_\Gamma d&=-\Re\frac{1}{2\pi\ii m_\Gamma}
                     \oint_\Gamma\operatorname{tr}H_1(s,d)\,ds,\\
 \boxed{|\mathcal R_\Gamma(d)|\le\beta_\Gamma
 :=\frac{\ell_\Gamma}{2\pi m_\Gamma}
       [e_\Gamma+r\varphi(\kappa_\Gamma)] .}
 \label{eq:regional-remainder}
\end{align}
\begin{proof}
The trace of $H-H_1$ is bounded by its nuclear norm.
Also $|\operatorname{tr}H^\ell|\le r\|H\|_2^\ell$,
including nonnormal $H$. Summing terms $\ell\ge2$ gives
$r\varphi(\kappa_\Gamma)$; contour integration proves the bound.
\end{proof}
A simpler conservative expression uses
$\varphi(x)\le x^2/[2(1-x)]$.
For the logarithm term, any consistently used induced matrix norm
can replace the Euclidean norm: each eigenvalue of $H$ is bounded
by that norm, hence $|\operatorname{tr}H^\ell|\le r\kappa^\ell$.
Pointwise verified bounds may instead be integrated along $\Gamma$.
Numerical errors in $g_\Gamma(p_0)$ and $J_\Gamma$ must be included:
if their certified errors are $e_0$ and $e_{J,a}$ and
$|d_a|\le r_a$, add $e_0+\sum_a e_{J,a}r_a$ to $\beta_\Gamma$.
No eigenvalue Hessian or fixed modal pattern is required.

\subsection{The actual replacement contract gives a small factorization}
\status{EXACT\_IDENTITY}
For the repository's exported fixed-support, fixed-equilibrium sharing
chart, define
\begin{align*}
 W_\rho&=\diag(1-\rho_1,1-\rho_1,\ldots,1-\rho_n,1-\rho_n),\\
 G(\rho)&=Y+W_\rho D_s,&
 N(\rho)&=W_\rho C_s+(I-W_\rho)C_f,\\
 V(\rho)&=-G(\rho)^{-1}N(\rho),&
 C(\rho)&=E_\theta+H_vV(\rho),\\
 B(K)&=B_p\diag K_p+B_i\diag K_i,&
 E_\tau(s)&=\diag(e^{-s\tau_i}).
\end{align*}
Here $V$ is the algebraic voltage \emph{Jacobian}, not the voltage phasor.
With ordinary gain coordinates,
\[
 \Delta(s,p)=sI-A_{\rm dev}-B_vV(\rho)-B(K)E_\tau(s)C(\rho).
\]
This is the inspected parametric model, not a universal identity for
all SG/GFL models or dispatch changes. If equilibrium/device matrices
change under another contract, their derivatives and remainders must
also be bounded. Endpoint state removal requires a separate chart.

Denote reference matrices by a subscript $0$, and set
\[
 J=G_0^{-1}\delta G,\quad
 V_1=-G_0^{-1}(\delta N+\delta G V_0),\quad
 \delta V=(I+J)^{-1}V_1,\quad
 V_{\rm rem}=- (I+J)^{-1}J V_1 .
\]
Thus $\delta V=V_1+V_{\rm rem}$ exactly. If uniformly
$\|J\|_2\le a<1$ and $\|V_1\|_2\le v$, then
\begin{equation}
 \|\delta V\|_2\le\frac{v}{1-a},\qquad
 \|V_{\rm rem}\|_2\le\frac{av}{1-a}. \label{eq:voltage-remainder}
\end{equation}
Because $G,N$ are affine in $\rho$, computable bounds are
\begin{align*}
a&=\sum_j r_{\rho,j}\|G_0^{-1}G_{\rho_j}\|_2,\\
v&=\sum_j r_{\rho,j}\|G_0^{-1}(N_{\rho_j}+G_{\rho_j}V_0)\|_2.
\end{align*}
These require verified solves and norm enclosures for a certificate.

Set $D_0(s)=B_v+B_0E_\tau(s)H_v$ and define
\begin{align}
 \mathcal U(s)&=[\,D_0(s)\quad B_pE_\tau(s)\quad B_iE_\tau(s)\,],\\
 \mathcal V(s,d)&=-
 \begin{bmatrix}
 \delta V\\
 \diag(\delta K_p)(C_0+H_v\delta V)\\
 \diag(\delta K_i)(C_0+H_v\delta V)
 \end{bmatrix}. \label{eq:physical-lowrank}
\end{align}
Direct expansion proves $\Delta-\Delta_0=\mathcal U\mathcal V$.
The linear part uses $V_1,C_0$ in the three block rows.
The remainder uses respectively
$V_{\rm rem}$, $\diag(\delta K_p)H_v\delta V$, and
$\diag(\delta K_i)H_v\delta V$, with the same minus signs.
For this export, $\mathcal U$ has 40 columns for 204 states:
20 voltage coordinates and two sets of ten PLL gain channels.
The dimension is not a discarded-state approximation.

For explicit conservative constants, let
$v_b=v/(1-a)$, $v_r=av/(1-a)$,
$k_p=\max_i|\delta K_{p,i}|_{\max}$,
$k_i=\max_i|\delta K_{i,i}|_{\max}$, and
$c_b=\|C_0\|_2+\|H_v\|_2v_b$.
Then
\begin{align*}
 \|\mathcal V\|_2&\le v_{\mathcal V}
   :=\sqrt{v_b^2+(k_pc_b)^2+(k_ic_b)^2},\\
 \|\mathcal V_{\rm rem}\|_2&\le v_{\rm rem}
   :=\sqrt{v_r^2+(k_p\|H_v\|_2v_b)^2+(k_i\|H_v\|_2v_b)^2}.
\end{align*}
If $w_\Gamma$ uniformly bounds
$\|\Delta_0(s)^{-1}\mathcal U(s)\|_2$ on the entire contour, use
$\kappa_\Gamma=v_{\mathcal V}w_\Gamma$ and
$e_\Gamma=r v_{\rm rem}w_\Gamma$.
Fixed nonsingular state/port scalings may improve these bounds if
applied consistently. A large bound is inconclusive, not instability.

\subsection{Affine descriptor realization with no parameter remainder}
\label{sec:affinelift}
\status{EXACT\_IDENTITY}
The inverse-network remainder can be avoided altogether for this
sharing contract. Retain the algebraic voltage perturbation $v$ and
the detector perturbation $e$ with the dynamic state $x$:
\begin{equation}
 \mathcal F(s,p)=
 \begin{bmatrix}
 sI-A_{\rm dev}&-B_v&-B(K)E_\tau(s)\\
 N(\rho)&G(\rho)&0\\
 -E_\theta&-H_v&I
 \end{bmatrix}. \label{eq:affineDAE}
\end{equation}
These are additional algebraic coordinates, not invented dynamics
or additional control channels. Eliminating $e$ and $v$ proves
\begin{equation}
 \det\mathcal F(s,p)=\det G(\rho)\det\Delta(s,p).
 \label{eq:descriptor-equivalence}
\end{equation}
On a region where $G$ is nonsingular, the finite roots and their
multiplicities coincide. The factor $\det G$ is independent of $s$.

Define $d=(\delta\rho,\delta K_p,\delta K_i)$ and the repeated diagonal
\[
 \Theta(d)=\diag(\delta\rho_1,\delta\rho_1,\ldots,
                \delta\rho_n,\delta\rho_n,\delta K_p,\delta K_i).
\]
With the reference operator $\mathcal F_0$, set
\begin{align}
 \mathcal U&=
 \begin{bmatrix}
 0&B_p&B_i\\ I&0&0\\0&0&0
 \end{bmatrix},&
 \mathcal V(s)&=-
 \begin{bmatrix}
 C_s-C_f&D_s&0\\
 0&0&E_\tau(s)\\
 0&0&E_\tau(s)
 \end{bmatrix},\\
 T(s)&=\mathcal V(s)\mathcal F_0(s)^{-1}\mathcal U,&
 H(s,d)&=\Theta(d)T(s). \label{eq:affine-return}
\end{align}
Then, exactly,
\[
 \mathcal F(s,p)-\mathcal F_0(s)
       =\mathcal U\Theta(d)\mathcal V(s),\qquad
 \frac{\det\mathcal F(s,p)}{\det\mathcal F_0(s)}
       =\det(I+\Theta(d)T(s)).
\]
The descriptor has dimension 234 and the determinant update has
dimension 40 for this export. In particular,
\[
 \boxed{H=H_1,\qquad H_{\rm rem}=0.}
\]
All nonlinear dependence of the finite root change is now in the
logarithm. The ordinary gain coordinates matter: log-gain increments
would require exponential parameter remainders.
Applying \eqref{eq:regional-log} to $\mathcal F$ gives the same
physical root means, and \eqref{eq:regional-remainder} simplifies to
\begin{equation}
 \boxed{\beta_\Gamma
   =\frac{\ell_\Gamma r}{2\pi m_\Gamma}
          \varphi(\kappa_\Gamma).} \label{eq:affine-remainder}
\end{equation}
The baseline-moment and coefficient integration errors still need
their separate enclosures.

\paragraph{Use the declared parameter box, not a product of large norms.}
For increments $l_a\le d_a\le u_a$, put
$r_a=\max(|l_a|,|u_a|)$. Let $R_d$ contain these nonnegative
displacement radii, repeated for the two occurrences of each
replacement coordinate. They equal box half-widths only when
$p_0$ is its center. Entrywise,
\[
 |H(s,d)|\le M(s):=R_d|T(s)|.
\]
For any fixed vector $w>0$ on a contour panel,
\begin{equation}
 \|H(s,d)\|_{w,\infty}
 \le \max_a\frac{(M(s)w)_a}{w_a}
 =:\kappa_w(s). \label{eq:positive-envelope}
\end{equation}
This follows by scaling with $\diag w$ and taking absolute row sums.
It is a standard positive-matrix bound, not a new stability theorem.
The repeated-coordinate correlation is relaxed conservatively;
the controllers remain local and the physical gain bounds unchanged.
A Perron-based weight can suggest $w$, but certification uses verified
inequalities for the chosen positive vector, not a floating-point
eigenvalue treated as an exact bound.

\paragraph{A concrete continuous-panel enclosure.}
On a panel contained in $|s-s_c|\le h$, put
$Q_c=\mathcal F_0(s_c)^{-1}$,
$W_c=Q_c\mathcal U$, and
\[
 e_\tau(h)=\max_i e^{-\Re(s_c)\tau_i}
                     (e^{h\tau_i}-1).
\]
For fixed delays,
\[
 \|\mathcal F_0(s)-\mathcal F_0(s_c)\|_2
 \le h+\|B_0\|_2e_\tau(h),\qquad
 \|\mathcal V(s)-\mathcal V(s_c)\|_2\le\sqrt2\,e_\tau(h).
\]
If verified bounds give
$a_c:=\|Q_c\|_2[h+\|B_0\|_2e_\tau(h)]<1$,
the Neumann series yields
\begin{equation}
 \|T(s)-T(s_c)\|_2
 \le
 \frac{\|\mathcal V(s_c)\|_2 a_c+\sqrt2\,e_\tau(h)}
      {1-a_c}\,\|W_c\|_2=:\epsilon_{T,c}. \label{eq:panel}
\end{equation}
Hence $R_d(|T(s_c)|+\epsilon_{T,c}\one\one^T)$ is an
entrywise panel envelope. Verified rounding/solve errors in
$Q_c,W_c,T(s_c)$ must be added; these formulas do not make an
ordinary dense solve self-certifying.
One positive-vector test on that envelope bounds the entire panel
and the entire parameter box. A finite covering establishes this contour condition. The displayed constants
apply in these coordinates; any nonunit state/equation scaling requires
rederiving the perturbation bounds in the scaled coordinates.
Thus the outstanding computation has an explicit error formula,
rather than an unspecified Hessian-sampling step.

\subsection{Regional maximum and a controller-eliminated witness}
\status{PROVED\_REDUCED\_MODEL}
Stack the necessary cluster inequalities and their verified errors.
Let $d=(\delta\rho,\delta K)$ range over compact boxes
$\mathcal R\times\mathcal K$, intersected with the original physical
limits. Let $A=J_\rho$, $C=J_K$, $b=-g_0+\beta$.
Every fully secure design in this region necessarily satisfies
\[
 A\delta\rho+C\delta K\le b.
\]
Consequently the LP
\begin{equation}
 U_{\mathcal D}=\max_{\delta\rho\in\mathcal R,\,
                         \delta K\in\mathcal K}
 (P^0)^T\delta\rho\quad
 \text{subject to }A\delta\rho+C\delta K\le b \label{eq:regionalLP}
\end{equation}
gives
\begin{equation}
 \boxed{P_F\le P^\star_{\mathcal D}
 \le (P^0)^T\rho_0+U_{\mathcal D}=P_{U,\mathcal D}.}
 \label{eq:regional-bracket}
\end{equation}
The lower candidate must belong to this same region and pass the full
spectrum and all declared events. Numerically validated feasibility
remains numerical evidence unless its own errors are certified.
Omitting frequency, RoCoF or actuator constraints from this LP is an
optimistic relaxation, not a safety claim.

For any $\eta\ge0$, the support functions of the two boxes give
\begin{equation}
 U_{\mathcal D}\le
 \eta^Tb+h_{\mathcal R}(P^0-A^T\eta)
          +h_{\mathcal K}(-C^T\eta). \label{eq:regionaldual}
\end{equation}
Minimization over $\eta$ recovers the feasible bounded LP value
by ordinary LP duality. For an interval box,
$h(z)=\sum_j\max(l_jz_j,h_jz_j)$.
The bound is valid for \emph{all} replacement directions in the
declared region, unlike the earlier fixed-ray calculation.

For a proposed replacement $\delta\rho$, the strict inequality
\begin{equation}
 \eta^T(A\delta\rho-b)>h_{\mathcal K}(-C^T\eta)
 \label{eq:regional-exclusion}
\end{equation}
proves that no permitted local gain vector in the region can satisfy
the necessary conditions, hence none can meet the full security contract.
The weights identify which collective requirements conflict.
If $C^T\eta=0$, their first-order aggregate has no gain authority;
finite residual authority is still allowed by $\eta^T\beta$.
Only the strict finite inequality, not a small derivative alone,
establishes exclusion. This is a regional statement, not a global
impossibility theorem for all local controllers.

\subsection{Where graph interactions actually enter}
\status{EXACT\_IDENTITY}
Group the rows/columns of $H$ by intervention site (two algebraic
voltage-equation channels and two gain channels per candidate).
The exact series contains
\[
 \operatorname{tr}H^2
 =\sum_i\operatorname{tr}H_{ii}^2
   +\sum_{i\ne j}\operatorname{tr}(H_{ij}H_{ji}),
\]
and $\operatorname{tr}H^\ell$ sums closed walks of length $\ell$
through these dynamical intervention ports. Their weights include
the full resolvent, P/Q coupling, SG dynamics and exact delay phases.
They are not necessarily physical transmission-line cycles, and $H$
is neither a Laplacian nor a physical damping matrix.
Even its diagonal blocks depend on the full network.

Writing $H=D+O$ with $D$ block diagonal gives
\[
 \det(I+H)=\det(I+D)\det(I+T),\qquad T=(I+D)^{-1}O,\quad
 \operatorname{tr}T=0.
\]
When the relevant logarithm series converge, the first omitted
inter-site term after discarding $O$ is
$-\tfrac12\operatorname{tr}T^2$ in the log determinant.
Its contribution to the root-sum change has the opposite sign
because of \eqref{eq:regional-log}; it need not be stabilizing.
If $\|T\|_2\le t<1$, the discarded log-determinant magnitude is
at most $r\varphi(t)$.
This supplies a falsifiable approximation error for neglecting
collective interactions, rather than assuming graph roughness
predicts damping.

An orthogonal graph basis can reorganize these operators.
It preserves traces and Euclidean norms for the corresponding
similarity, but its diagonal truncation changes the approximation.
Low graph frequency alone does not justify that truncation.
No GSP advantage or physical cycle mechanism is established until
the retained/off-diagonal terms explain a measured limiting decision.

\section{Collective curvature and a graph-resolved finite remainder}
\label{sec:walkbounds}
\status{EXACT\_IDENTITY}
Let $d\in\mathbb R^{3n}$ denote ordinary replacement and gain increments.
Define fixed diagonal selectors $E_a$ such that
$\Theta(d)=\sum_a d_a E_a$: a replacement selector contains two unit
entries, whereas a gain selector contains one. In the affine descriptor
chart, $H(s,d)=\sum_a d_a E_aT(s)$ exactly.
All delays remain fixed inside $T(s)$.

\begin{theorem}[Collective cluster curvature]
Under the contour hypotheses of Section~\ref{sec:regional}, the real
cluster margin has the second-order expansion
\begin{equation}
 g_\Gamma(p_0+d)=g_{\Gamma,0}+J_\Gamma d+
       \tfrac12d^TQ_\Gamma d+\mathcal R_{\Gamma,3}(d),\qquad
 (Q_\Gamma)_{ab}=
 \Re\frac{1}{2\pi\ii m_\Gamma}
       \oint_\Gamma\operatorname{tr}(E_aTE_bT)\,ds .
 \label{eq:collective-hessian}
\end{equation}
The matrix $Q_\Gamma$ is real symmetric. It is not required to be
positive semidefinite and is not a Laplacian.
\end{theorem}
\begin{proof}
Substitute $H=\sum_a d_aE_aT$ in the convergent series for
$-\operatorname{tr}\log(I+H)=-\operatorname{tr}H+
\operatorname{tr}H^2/2-\cdots$ in \eqref{eq:regional-log}.
Cyclic invariance of the trace gives
$\operatorname{tr}(E_aTE_bT)=\operatorname{tr}(E_bTE_aT)$.
Collecting terms proves the formula. Uniform convergence on the
compact contour permits termwise integration. This differentiates a
cluster moment, so internal root collisions do not require eigenvector
derivatives or a simple-root assumption.
\end{proof}
For scalar selectors the integrand is $T_{ab}T_{ba}$.
Thus a mixed change at two sites follows a return path through the
full dynamical network. Two individually favorable first-order
changes need not remain favorable when combined. The sign is set by
the complex contour integral, not by graph adjacency or a positive
edge weight. Diagonal terms already include the full grid.

\status{PROVED\_REDUCED\_MODEL}
Suppose a continuous, piecewise bounded nonnegative envelope satisfies
\[
 |H(s,d)|\le M(s),\qquad
 \sup_{s\in\Gamma}\rho(M(s))<1,\qquad d\in\mathcal D_d .
\]
A verified positive-vector inequality on every panel is one sufficient
way to establish this condition. Set
\[
 \Psi_k(M)=\sum_{\ell=k}^{\infty}\frac{\operatorname{tr}M^\ell}{\ell}
 =-\log\det(I-M)-
     \sum_{\ell=1}^{k-1}\frac{\operatorname{tr}M^\ell}{\ell},
 \qquad k\ge2 .
\]
Here the determinant is positive and the logarithm is real: real
eigenvalues of $M$ are below one, and nonreal eigenvalues occur in
conjugate pairs. Despite that spectral representation, every term in
the defining trace series is nonnegative.

\begin{theorem}[Closed-walk remainder and interaction-omission budget]
The remainder after degree $k-1$ obeys
\begin{equation}
 \boxed{|\mathcal R_{\Gamma,k}(d)|\le
 \beta_{\Gamma,k}^{\rm walk}:=
 \frac{1}{2\pi m_\Gamma}
       \oint_\Gamma\Psi_k(M(s))\,|ds|.}
 \label{eq:walk-remainder}
\end{equation}
Partition channels into physical intervention sites and let
$H_{\rm b}=\operatorname{blockdiag}(H_{11},\ldots,H_{nn})$.
The omitted contribution to the cluster-shift formula, obtained by
replacing $H$ with $H_{\rm b}$, is bounded by
\begin{equation}
 \boxed{\beta_\Gamma^{\rm cross}=
 \frac{1}{2\pi m_\Gamma}\oint_\Gamma
 \left[\Psi_2(M)-\sum_i\Psi_2(M_{ii})\right]|ds|.}
 \label{eq:cross-budget}
\end{equation}
This is a bound on an algebraic approximation to the contour formula;
it does not assert that $H_{\rm b}$ is a realizable physical grid.
\end{theorem}
\begin{proof}
The entrywise inequality $|H^\ell|\le M^\ell$ implies
$|\operatorname{tr}H^\ell|\le\operatorname{tr}M^\ell$.
Sum the tail of the logarithm series and integrate its absolute value
to obtain \eqref{eq:walk-remainder}.
For \eqref{eq:cross-budget}, the terms remaining after block-diagonal
subtraction are exactly index sequences that visit more than one
block before returning to their starting channel. Bounding the
absolute value of each product by the corresponding product of
entries of $M$ gives
$\operatorname{tr}M^\ell-\sum_i\operatorname{tr}M_{ii}^\ell$.
The degree-one difference is zero. Summation proves the result.
\end{proof}

The earlier scalar bound follows from
$\operatorname{tr}M^\ell\le r\kappa^\ell$ whenever a common induced
norm of $M$ is bounded by $\kappa<1$. Retaining the actual walk sums
can therefore substantially reduce conservatism. A computable,
slightly looser expression avoids logarithms:
\[
 \Psi_k(M)\le\frac1k\operatorname{tr}\!\left[M^k(I-M)^{-1}\right].
\]
All inverse, determinant, integration and rounding errors must still
be enclosed. In particular, a verified evaluation of
\eqref{eq:cross-budget} requires an \emph{upper} enclosure of the
total minus \emph{lower} enclosures of the within-site terms;
subtracting two upper bounds would be invalid. Directly summing
cross-site walks is an alternative.

\paragraph{Meaning of beyond nodal damping.}
Equation~\eqref{eq:collective-hessian} identifies the signed,
frequency-dependent interaction between two replacement or tuning
actions. Equation~\eqref{eq:cross-budget} quantifies the error made
by dropping those interactions. These are not nodal damping
constants: $T$ contains the grid resolvent, SG dynamics, voltage/PQ
coupling and the exact phases $e^{-s\tau_i}$. They supplement the
physical torque--speed damping operator derived in the appendices;
they do not turn a contour return matrix into a physical impedance.
The closed walks are paths in the intervention-port representation,
not automatically cycles of transmission lines. No communication
between PLLs has been introduced.

\subsection{A stronger necessary replacement relaxation}
Write $d\in[l,u]$ and introduce a real symmetric matrix $Z$.
Every actual design admits $Z=dd^T$, hence satisfies
\begin{align}
 &\begin{bmatrix}1&d^T\\d&Z\end{bmatrix}\succeq0,\\
 &Z_{ab}\ge l_ad_b+l_bd_a-l_al_b,\qquad
 Z_{ab}\ge u_ad_b+u_bd_a-u_au_b,\\
 &Z_{ab}\le u_ad_b+l_bd_a-u_al_b,\qquad
 Z_{ab}\le l_ad_b+u_bd_a-l_au_b .
\end{align}
Retain the diagonal case of these inequalities as well.
Impose, for each verified cluster,
\begin{equation}
 g_{\Gamma,0}+J_\Gamma d+
       \tfrac12\langle Q_\Gamma,Z\rangle
       \le\beta_{\Gamma,3}^{\rm walk}.
 \label{eq:quadratic-outer}
\end{equation}
Maximizing $(P^0)^T\delta\rho$ over this semidefinite relaxation gives
an optimistic regional replacement upper bound. Dropping
$Z=dd^T$ enlarges the feasible set, which is the required direction;
the returned gains need not be physically feasible.
Neither a primal floating-point optimum nor a tiny solver residual
is a certified upper bound. A feasible dual with enclosed errors,
or another verified objective bound, is required for that label.

Certified coefficient errors add
$e_0+\sum_a e_{J,a}r_a+
\tfrac12\sum_{a,b}e_{Q,ab}r_ar_b$ to the right-hand side,
where $r_a=\max(|l_a|,|u_a|)$. The lift may sharpen the linear
relaxation, but no universal dominance follows if different
remainders or additional relaxations are used. Both outer problems
can be intersected because each contains all secure designs.
This strengthening has been derived and identity-tested; the
IEEE--39 semidefinite upper-bound problem has not been solved.

\section{A reproducible analytical design and bound procedure}
\label{sec:algorithm}
The optimization and certification roles are distinct. The exact
all-PLL formula in Eq.~\eqref{eq:allgain} constructs candidates.
The regional outer problem permits arbitrary admissible modal
patterns and limits what any such candidate can achieve.

\begin{enumerate}
\item Freeze the sharing convention, fixed delay vector, physical
gain limits, replacement region, event set, and output definitions.
Solve and audit the equilibrium; verify the algebraic chart.
\item At a candidate center, assemble the exact descriptor with
exponential delay factors. Use numerical roots to propose contours;
verify their counts, boundary regularity and parameter-box envelopes.
Numerical root locations alone do not close this step.
\item Integrate the reference moments and $J_\Gamma$ (and, if used,
$Q_\Gamma$) with error enclosures. Evaluate the walk remainder.
If a panel or algebraic chart cannot be verified, subdivide it or
return an unresolved box; do not fill it by interpolation.
\item Solve the necessary LP or lifted outer problem and verify an
upper objective bound or a strict controller-exclusion witness.
Compare this upper bound against the box-only replacement ceiling.
\item Use the analytical gain map or a trust-region local step to
propose a design, then evaluate the full delayed spectrum and all
events. Failed designs do not supply a feasible lower bound.
\item Report the same-domain feasible and upper values separately,
with their uncertainty. Subdivide parameter regions only under a
declared branch rule if a wider bound is sought.
\end{enumerate}
No global convergence theorem for the predictor/corrector is claimed.
A global upper bound would require a verified cover of the full
design domain, including all support charts and unresolved boxes.
The full event-feasible problem generally has more active constraints
than the two-row reduced LP; its SG-retention structure cannot be
inferred from a two-anchor theorem.
\section{Targeted validation of the all-PLL construction}\label{sec:mapvalidation}
\status{NUMERICALLY\_VALIDATED}
This subsequent experiment tests Eq.~\eqref{eq:allgain} with one
preregistered replacement step. It is not an optimization campaign.
All ten fractions change from $\rho_i=0.875$ to $0.876$ with
$\tau_i=40$ ms fixed. The target is the leading PLL-family pole
in the frozen 2--12 Hz catalog, not the globally rightmost pole:
\[
 \lambda_\star=-0.80156077+\ii\,30.85047169\ \mathrm{s}^{-1}.
\]
Its PLL-angle pattern is held fixed and all twenty gains are calculated
from the full-grid return. The trial design and modal predictions
were hashed before independent full-model validation.

Known gains were reconstructed for all eleven nonreal catalog roots,
with maximum relative error $7.266\times10^{-8}$. Nine vector derivative tests,
using three modes and buses 30, 34 and 38, had maximum relative error
$2.629\times10^{-5}$ and maximum absolute infinity error $2.965\times10^{-1}$.
The largest hidden-block condition number was $1.637\times10^{9}$;
these residuals establish numerical consistency, not physical accuracy
to all the reported digits.
Direct automatic differentiation of the nonlinear equations gave
maximum relative Jacobian discrepancy $3.682\times10^{-18}$ and maximum equilibrium
infinity residual $2.807\times10^{-9}$. This cross-check uses the same
nonlinear device equations; it is not an independent-simulator or
physical-data validation.

An independently assembled Jacobian of the same model, followed by the delayed calculation, preserved the target
pole within $1.511\times10^{-13}\ \mathrm{s}^{-1}$ with PLL-pattern MAC $1.000000000000$.
For the other modes, the largest relative finite-step prediction error
was 2.791\%; the preregistered 1\% criterion concerns derivative
validation, not a finite-step Taylor remainder bound.
A numerical argument-principle count, cross-checked against determinant
phase and closed by a tail bound after rotational deflation, found no
roots with $\Re s>-0.05\ \mathrm{s}^{-1}$ for the baseline, retuned
design, or fixed-gain comparator. Exact exponentials were used;
floating-point contour counts are not interval certificates.

The retuned design passed all five frozen nonlinear events, simulated
with the canonical method of steps (Rodas5P, tolerance $10^{-9}$,
maximum step $0.01$ s, horizon $60$ s). The 39-bus phase metrics use
$0.5$ s windows. Frequency and RoCoF limits are $0.5$ Hz and
$0.5$ Hz/s; normalized SG actuator slack must exceed $0.002$.
\begin{center}\small
\begin{tabular}{rrrrr}\toprule
Bus & Load change [MW] & $|\Delta f|$ [Hz] & RoCoF [Hz/s] & SG slack\\
\midrule
8 & -100 & 0.465412 & 0.138451 & 0.096623\\
16 & +100 & 0.447157 & 0.126491 & 0.006341\\
16 & -100 & 0.474624 & 0.129044 & 0.093434\\
29 & +100 & 0.431226 & 0.128104 & 0.011190\\
29 & -100 & 0.444619 & 0.200774 & 0.102894\\
\bottomrule\end{tabular}
\end{center}
All voltage samples remained within $[0.972062,1.078775]$ pu
against the declared $[0.9,1.1]$ limits.
The candidate contains 4732.818715 MW of GFL (87.600\%) and
669.942375 MW of retained SG: a prescribed increase of 5.402761 MW.
The full vectors are in
\path{experiments/all_pll_gain_map_validation_20261003/designs/analytic.toml}.

\paragraph{Interpretation and negative evidence.}
This closes the numerical chain from an analytical gain construction
to a feasible full-model candidate for the declared events.
It does not establish a maximum or demonstrate that retuning is needed
for this step: the same replacement with unchanged gains also passes
the spectral margin. The subsequent matched nonlinear comparison
in Section~\ref{sec:posterclosure} finds that both designs pass all five
events, with small opposing changes in frequency and RoCoF metrics.
Neither superiority nor exact performance equivalence is established.
Two frozen eigenpairs/patterns produce nonzero assignment residuals
after replacement, but this does not prove that local PLLs cannot
stabilize the network when their modal patterns are allowed to change.
No heterogeneous-delay, neighboring-feedback, GSP-compression,
convergence, or global-optimality claim follows from this trial.
Data, prediction hashes, code, trajectories and claim limits are in
\path{experiments/all_pll_gain_map_validation_20261003/}.
% END GENERATED ALL-PLL VALIDATION


\paragraph{Reproducible gain vector.}
For buses 30--39, every listed replacement fraction is $\rho_i=0.876$.
These are constructed feasible-candidate gains, not optimal gains.
\begin{center}\small
\begin{tabular}{rrr}\toprule
Bus & $K_{p,i}$ & $K_{i,i}$\\\midrule
30 & 28.27608095 & 246.76647162\\
31 & 28.24652453 & 245.80547706\\
32 & 28.24503605 & 245.70201441\\
33 & 28.21712688 & 244.76013817\\
34 & 28.17474906 & 246.56932577\\
35 & 28.24738622 & 245.39990024\\
36 & 28.25363221 & 246.19628049\\
37 & 28.27912035 & 246.57391812\\
38 & 28.25524609 & 245.98556358\\
39 & 28.28100338 & 246.76474736\\
\bottomrule\end{tabular}\end{center}

\section{Bound checks, falsification and regional evidence}
\label{sec:boundchecks}
\status{NUMERICALLY\_VALIDATED}
The descriptor determinant and affine-factor identities were checked
at the three frozen frequencies 0.5, 5 and 10 Hz, with
$s=-0.05+\ii 2\pi f$, heterogeneous delays uniformly spread from
0 to 40 ms, replacement displacement radii 0.001, and ordinary gain
radii 1\%. This delay pattern is different from the uniform-delay
nonlinear trial; their lower and upper evidence cannot be mixed.
The largest previously recorded normalized identity discrepancy was
$1.011\times10^{-12}$.

The new checks retain those same points and box. They compare the
scalar logarithm-tail bound with the entrywise closed-walk majorant.
The table contains pointwise bounds on the trace logarithm, not
eigenvalue shifts or a sampled replacement certificate.
\begin{center}\small
\begin{tabular}{rrrr}\toprule
$f$ [Hz] & scalar tail & walk tail & reduction factor\\\midrule
0.5 & 0.0189751 & 0.000964533 & 19.67\\
5 & 1.38717 & 0.0303027 & 45.78\\
10 & 0.0056687 & 0.000313074 & 18.11\\
\bottomrule\end{tabular}\end{center}
The first- and second-order trace predictions and the cross-site
omission budgets passed all three pointwise numerical checks.
This supports improved bound sharpness at these points; it does
not demonstrate that graph terms increase attainable replacement.

For the separate two-state pure-DDE fixture
\[
 \Delta_{\rm toy}(s,p)=
 \diag(s+1,s+2)-pe^{-s/10}
 \begin{bmatrix}1&1\\1&-1\end{bmatrix},\qquad |p|\le 0.01 ,
\]
the circle centered at $-1$ with radius $1/4$ admits the analytic
whole-circle bound $\kappa\le12/175<1$ and one enclosed root.
Implicit differentiation gives
$\lambda'(0)=e^{1/10}$ and $\lambda''(0)=1.8e^{1/5}$.
The contour Hessian reproduces both values to numerical error.
For the five frozen values of $p$ at 256 and 512 quadrature nodes,
the largest quadratic error is $3.140\times10^{-6}$, below the
analytic cubic-tail bound $5.770\times10^{-5}$.
This is a theorem/check on the declared fixture, not on the power grid.

\subsection{Continuous IEEE--39 contour test}
Complex ball arithmetic \cite{arb} at 128-bit precision encloses every point
of a square centered at the preserved slow-mode estimate
$-0.12229688588491999+\ii\,0.32408378548166134$, with half-side
0.04. The region uses uniform fixed delay 40 ms,
$\rho_i\in[0.874,0.876]$ and ordinary gains within 1\% of the
baseline. The matrix inputs are the exact binary64 values parsed
from the frozen exports. The fixed delay is the exact decimal
$0.04=1/25$ s enclosed in ball arithmetic, not its binary64 rounding.
Physical and equilibrium uncertainty
are not enclosed by this numerical certificate.
An interval inverse verifies $G(\rho)$ throughout that box.
The contour cover uses 485 accepted panels and gives
$\kappa\le 0.98794235<1$.
The enclosed argument-principle integral contains exactly one
integer, 1. Consequently one root remains inside this square for
every permitted parameter combination. Its real part is below
$-0.08229\ \mathrm{s}^{-1}$, so this one cluster satisfies the
$-0.05\ \mathrm{s}^{-1}$ margin throughout the box.
This verifies one cluster; it is not a complete-spectrum or nonlinear
certificate.
All ball enclosures, parameters and code are saved in the
reproduction package. The coarse first-order scalar moment remainder
is bounded by $4.458179\ \mathrm{s}^{-1}$, much larger than the
required margin. This bound is conservative even though the contour
geometry proves this cluster safe. The sharper walk majorant has
not yet been integrated with verified errors around this contour.

\subsection{A decisive negative result about the chosen regional box}
\status{NEGATIVE\_RESULT}
To the reported numerical precision, the feasible 87.6\% design has
$\rho_i=0.876$ at every candidate and its gains lie inside the
uniform-delay verification box. Its reported objective coincides with
the box-only upper bound
\[
 P_{\rm box}=0.876(P^0)^T\mathbf1=4732.818715\ {\rm MW}.
\]
Conditional on true feasibility of the exact upper-corner design,
this is the optimum \emph{of that box}, by monotonicity of a positive
dispatch objective. No spectral theorem is needed for this fact.
It is not an irreducible SG-support bound and is not a meaningful
near-global-optimality result. A purported upper bound below this
value would contradict the numerical witness and require resolving
that discrepancy. For strict arithmetic bookkeeping, the binary64
literal $0.876$ exceeds the exact decimal endpoint by
$1/1125899906842624000$; the numerical witness is not treated as
an interval-certified corner point. No rigorous lower--upper gap
is claimed from these rounded values.
Thus even perfect verification of this regional box cannot establish
the main application claim of a physical interior replacement limit.
The manuscript does not expand the box after observing this result.





\section{Targeted closure: physical damping and matched nonlinear comparison}
\label{sec:posterclosure}
This follow-up adds two declared tests without changing the historical
design, events, gains, or frozen result files. Its protocol and outputs
are stored in \texttt{experiments/poster\_julia\_validation\_20261003/}.
The harmonic tests are algebraic; the event comparator uses Julia 1.11.9.

\subsection{Finite-delay physical damping law in IEEE--39}
\status{NUMERICALLY\_VALIDATED}
For the historical baseline and analytical designs, at each of
$f=0.5,5,10$ Hz, increase the delay at one of buses 30--39 by exactly
the numerical increment 1 ms, starting from uniform 40 ms.
All $2\times3\times10=60$ cases are retained. No case was chosen by
its sign or stability outcome. At each point, independently eliminate
the 194 hidden states in the old and modified 204-state operators and
compare the difference with the finite rank-one formula
in Appendix~\ref{app:physical}.
The physical mass scaling is
$M_i=2H_iS_i^0(1-\rho_i)$, in MW\,s for per-unit speed.

The largest relative error in $\delta Z=\gamma_i u_i v_i^*$ is
$4.85\times10^{-8}$. Every computed Hermitian increment has one
positive and one negative eigenvalue above the declared relative
threshold $10^{-7}\|\delta D_H\|_2$; the other eight eigenvalues
have relative magnitude at most $1.40\times10^{-9}$.
The minimum squared sine between the two propagation paths is
$0.00327$, so none of these sampled paths is numerically collinear.
Hidden-block condition numbers and all signed eigenvalues are
reported in \texttt{TABLE\_02\_FINITE\_DELAY\_RANK\_LAW.csv}.
This supports the finite physical identity at the declared points.
It does not certify an interval of frequencies, imply stability from
the sign of the increment, or improve a replacement-capacity bound.

\begin{figure}[htbp]\centering
\includegraphics[width=.87\linewidth]{reports/poster/ias2026/collective_interaction_20261003/generated/figures/rank_two.pdf}
\caption{Physical damping redistribution at the preregistered illustrative
frequency 5 Hz, analytical design. One PLL delay is increased from
40 to 41 ms at a time. The symmetric-log scale exposes both signs.
This is a harmonic linear-model check, not a nonlinear stability test.}
\label{fig:rankclosure}
\end{figure}

\subsection{Same replacement, unchanged versus analytical gains}
\status{NUMERICALLY\_VALIDATED}
The missing matched comparator is now executed at the same
$\rho_i=0.876$ and $\tau_i=40$ ms as the analytical candidate.
Its gains remain at the common baseline values.
The canonical method-of-steps solver uses Rodas5P, tolerance $10^{-9}$,
maximum step 0.01 s, and a 60 s event horizon. Both phase-derived
frequency and RoCoF metrics use the canonical 0.5 s window over all
39 bus voltages. The five nominal event magnitudes correspond to
the existing constant-impedance perturbation contract.
Equilibrium residual and gain bounds are checked before execution.

\begin{table}[htbp]\centering\small
\caption{All five matched events. $F$ is peak absolute frequency deviation
in Hz; $R$ is peak absolute RoCoF in Hz/s. A denotes the analytical map;
fixed retains the original gains. All original guards also pass.}
\begin{tabular}{rrrrrl}\toprule
Bus, MW & $F_{\rm fixed}$ & $F_{\rm A}$ & $R_{\rm fixed}$ & $R_{\rm A}$ & Guards\\\midrule
8, -100 & 0.465209 & 0.465412 & 0.138636 & 0.138451 & Pass / Pass\\
16, +100 & 0.446979 & 0.447157 & 0.126583 & 0.126491 & Pass / Pass\\
16, -100 & 0.474417 & 0.474624 & 0.129118 & 0.129044 & Pass / Pass\\
29, +100 & 0.431049 & 0.431226 & 0.128359 & 0.128104 & Pass / Pass\\
29, -100 & 0.444418 & 0.444619 & 0.201209 & 0.200774 & Pass / Pass\\
\\[-2.5mm]\bottomrule
\end{tabular}\label{tab:pairedclosure}
\end{table}

Both designs pass all five events. Analytical retuning increases
frequency excursion by $0.000177$--$0.000207$ Hz and reduces RoCoF
by $0.0000744$--$0.000435$ Hz/s across these cases. Its SG actuator
slack is also slightly smaller. These are small numerical tradeoffs,
not a demonstrated engineering advantage; no refinement study was
performed to turn these small differences into accuracy-certified
performance claims. The purpose of the exact gain map is achieved:
the assigned PLL eigenpair is preserved after the prescribed
replacement step. This matched trial does not show that retuning
was necessary for feasibility, that it maximized replacement, or
that it dominates fixed gains.

\begin{figure}[htbp]\centering
\includegraphics[width=\linewidth]{reports/poster/ias2026/collective_interaction_20261003/generated/figures/events.pdf}
\caption{Preserved analytical-design trajectories for every frozen
event, plotted across the full 60 s horizon. Frequency and RoCoF use
the declared 0.5 s phase window. The new comparator's complete
trajectories and paired metrics are included separately in the archive.}
\label{fig:allfrozen}
\end{figure}

\paragraph{Updated contribution boundary.}
The physical finite-delay theorem, the collective contour Hessian,
and the walk remainder provide explicit equations for network-mediated
effects. The numerical evidence now includes all 60 physical-law
checks and a matched five-event comparator. No informative
verified physical replacement maximum follows. The posterior
comparison therefore supports a theory-and-validation poster;
the stronger application claim remains open.

\section{Discussion, contribution boundary and limitations}
\begin{center}\small
\begin{tabular}{p{0.29\linewidth}p{0.61\linewidth}}\toprule
Result & Evidentiary scope\\\midrule
Affine descriptor and gain map & Exact identities in the fixed chart.\\
Cluster Hessian and walk tails & Proved under the stated contour/envelope hypotheses.\\
Regional LP / lifted outer bound & Necessary relaxation; verified data and dual bound required.\\
87.6\% candidate, five events & Full-model numerical result for uniform 40 ms delay.\\
18--46-fold tail reduction & Three frozen algebraic points, not a capacity increase.\\
Physical maximum or global gap & Not established.\\
Local PLLs fundamentally insufficient & Not established with freely varying modal patterns.\\
GSP compression or physical cycle law & Not established.\\\bottomrule
\end{tabular}
\end{center}

\paragraph{What is technically added.}
The model-specific lift links replacement and PLL gain increments
to a finite characteristic perturbation without a network-inverse
parameter remainder. Cluster curvature resolves pairwise collective
actions, and the walk budget quantifies the error of dropping them.
The outer relaxation allows all admissible gains and modal patterns;
a strict finite dual witness would therefore be stronger than
failure to preserve a prescribed eigenvector. This is an applied
bound construction assembled from established tools. Bibliographic
priority for the combination has not been proved.

\paragraph{What still prevents a strong application paper.}
An informative verified regional upper bound or controller-exclusion
witness below the box-only ceiling is missing. The matched nonlinear fixed-gain comparison now shows both designs
feasible, with small opposing metric changes. A fair fixed-structure
spectral-abscissa/SQP comparator with identical constraints remains missing. A controller tuned with $I+L+L^2$ terms has not been
implemented or shown necessary. The graph Fourier basis may express
interactions, but no low-frequency truncation theorem follows.

\paragraph{Physical limits.}
The certificate concerns a local linearization of a specific
nonlinear model. Nonlinear event tests remain separate; no global
nonlinear stability, arbitrary-disturbance safety, or unmodeled
converter-current-limit guarantee follows. Delays may be
heterogeneous in the derivations, but the preserved nonlinear
feasible result has uniform delays. The graph moment in the
appendices concerns a declared swing-network class and has not
been transferred to the full IEEE--39 converter model.

\section{Conclusion}
Collective effects can be made explicit at the level of finite
parameter changes: the full delayed network produces a cluster
Hessian and a closed-walk remainder, which lead to necessary
regional replacement bounds. The derivations permit moving modes
and fixed heterogeneous delays without a low-frequency delay
approximation. The available IEEE--39 results validate algebraic
sharpness and one feasible analytical retuning construction.
One continuous complex-ball calculation also certifies retention and
margin stability of a slow cluster over its complete local box.
These results do not yet establish a physically limiting replacement ceiling
or superior co-design performance. The next decisive result is an
informative bound on a preregistered region extending beyond the
current feasible corner, followed by a comparison under the same
nonlinear security contract.

\section*{Reproducibility and research integrity}
The main source is edited in place; pre-edit source bytes and prior
experiment artifacts are preserved. New checks, protocol, exact
input copies, package versions, tables, claim ledger and manifests
are under \path{experiments/regional_paper_closure_20261003/}.
Numerical table entries in the new main text are generated from
the saved result files. The initial closure performed linear-algebra
and contour checks. The subsequent poster closure adds the 60 physical-law
checks and five new fixed-gain nonlinear comparator trajectories under
\path{experiments/poster_julia_validation_20261003/}.
The historical event evidence remains explicitly attributed to its
frozen experiment. Three independent reviews audited the equations,
evidence provenance and literature boundary; their corrections and
unresolved objections are recorded in the package.

\clearpage
\appendix
\section{First correction: derive the closed PLL before assigning damping}
\status{EXACT\_IDENTITY}
At a locked equilibrium $\theta_i^*=\arg V_i^*$, let
$V_i^\circ=|V_i^*|>0$ and let $\phi_i$ be terminal phase perturbation.
The linearized detector is $\delta e_i=V_i^\circ(\phi_i-\theta_i)$.
Eliminating the PLL frequency and integrator gives, for zero history,
\begin{equation}
 H_i(s)=\frac{\theta_i(s)}{\phi_i(s)}
 =\frac{V_i^\circ(K_{p,i}s+K_{i,i})e^{-s\tau_i}}
 {s^2(1+t_{f,i}s)+V_i^\circ(K_{p,i}s+K_{i,i})e^{-s\tau_i}}. \label{eq:H}
\end{equation}
For $K_{i,i}>0$, direct Taylor division proves
\begin{equation}
 H_i(s)=1-\frac{s^2}{V_i^\circ K_{i,i}}
 +\frac{K_{p,i}/K_{i,i}-t_{f,i}-\tau_i}{V_i^\circ K_{i,i}}s^3+O(s^4).
 \label{eq:series}
\end{equation}
Pure delay first appears at cubic order in this phase-tracking transfer,
not in its constant, linear or quadratic coefficients.
Any subsequent reduction using analytic factors and inverses nonsingular
at zero preserves the $O(s^3)$ delay difference. Dividing an angle
operator by $s$ shifts that difference to $O(s^2)$.
Singular elimination can change the order. The previously ill-conditioned
full hidden block at $s=0$ prevents an unqualified application of this
argument to the full reduced grid.

Consequently $L_\tau=\operatorname{Sym}(B_cK_iT_\tau C)$ is not established
for this implemented closed PLL. It may arise for another loop or
approximation, which would need its own derivation.
At finite frequency the inverse tracking defect is exactly
\begin{align}
 Q_i(s)&=H_i(s)^{-1}-1
 =\frac{s^2(1+t_{f,i}s)e^{s\tau_i}}{V_i^\circ(K_{p,i}s+K_{i,i})},\\
 |Q_i(\ii\Omega)|&=
 \frac{\Omega^2\sqrt{1+\Omega^2t_{f,i}^2}}
 {V_i^\circ\sqrt{K_{i,i}^2+\Omega^2K_{p,i}^2}},\\
 \arg Q_i(\ii\Omega)&=\pi+\arctan(\Omega t_{f,i})
 -\arctan(\Omega K_{p,i}/K_{i,i})+\Omega\tau_i\pmod{2\pi}.
\end{align}
Delay rotates this defect without changing its magnitude. This defect
is not itself physical damping or a Laplacian.

\section{Physical collective damping and its finite latency law}\label{app:physical}
\status{EXACT\_IDENTITY}
The exact characteristic operator of the full linearization is
\[
 \Delta(s)=sE-A_0-\sum_i b_ic_i^T e^{-s\tau_i}.
\]
The two delayed PLL injections share one column $b_i$, with entries
$K_{p,i}/t_{f,i}$ and $K_{i,i}$. $c_i^T$ differentiates the actual
detector after algebraic voltage elimination.

In exported $E=I$ coordinates, retain per-unit SG speed deviations $\nu$
and eliminate all other states $z$ wherever $\Delta_{zz}$ is invertible:
\begin{equation}
 Z(s)=M(\Delta_{\nu\nu}-\Delta_{\nu z}\Delta_{zz}^{-1}\Delta_{z\nu}),
 \qquad d=Z(s)\nu,\qquad D_H(\Omega)=\Herm Z(\ii\Omega). \label{eq:Z}
\end{equation}
$d_i$ is mechanical torque expressed as MW at nominal speed, injected
as $d_i/M_i$ in the swing equation. For physical harmonic phasors,
\[
 \langle d^T\nu\rangle=\tfrac12\widehat\nu^*D_H(\Omega)\widehat\nu .
\]
This is incremental mechanical supply, not total AC power or a nonlinear
Lyapunov function. A steady harmonic response needs to exist; an unstable
system may only admit a formal particular phasor solution.

\begin{theorem}[One delayed channel redistributes physical damping]
Fix $\rho$, equilibrium, gains, ports and $\Omega>0$.
Change only $\tau_i$ to $\tau_i+\delta\tau_i\geq0$.
Assume both hidden blocks are invertible. At $s=\ii\Omega$, evaluate
the following factors at the old design:
\begin{align*}
 \kappa_i&=e^{-s\tau_i}-e^{-s(\tau_i+\delta\tau_i)},&
 h_i&=c_{z,i}^T\Delta_{zz}^{-1}b_{z,i},\\
 u_i&=M(b_{\nu,i}-\Delta_{\nu z}\Delta_{zz}^{-1}b_{z,i}),&
 v_i^*&=c_{\nu,i}^T-c_{z,i}^T\Delta_{zz}^{-1}\Delta_{z\nu},\\
 \gamma_i&=\kappa_i/(1+\kappa_i h_i),&
 a_i&=\gamma_i u_i .
\end{align*}
Then
\begin{equation}
 \boxed{\delta Z=\gamma_i u_iv_i^*,\qquad
 \delta D_H=\tfrac12(a_iv_i^*+v_ia_i^*).} \label{eq:rank}
\end{equation}
If $\gamma_i\ne0$ and $u_i,v_i$ are independent, $\delta D_H$ has
exactly one positive and one negative eigenvalue, all others zero.
\end{theorem}
\begin{proof}
The full change is $\kappa_i b_ic_i^T$. Applying Sherman--Morrison
inside the Schur complement proves (\ref{eq:rank}). Nonsingularity of
both hidden blocks makes $1+\kappa_i h_i\ne0$ by the determinant lemma.
The range of $\delta D_H$ is contained in the span of $a_i,v_i$.
Its two possibly nonzero eigenvalues are
\[
 \ell_\pm=\frac{\Re(v_i^*a_i)\pm
 \sqrt{\|a_i\|^2\|v_i\|^2-[\Im(v_i^*a_i)]^2}}{2}.
\]
Their product is
\[
 \ell_+\ell_-=-\frac{|\gamma_i|^2}{4}
 \left(\|u_i\|^2\|v_i\|^2-|v_i^*u_i|^2\right)<0
\]
by strict Cauchy--Schwarz. This proves both signs and rank two.
\end{proof}
This is a finite-change statement for the full linearization, not
merely a first-order approximation. It says that a local latency
change increases harmonic dissipation in some collective speed
directions and decreases it in others, under the stated independence.
Hence its negative is not a positive-semidefinite delay penalty either.
Collinear paths can instead give a rank-one change of either sign or
zero; phase-periodic increments can make $\kappa_i=0$.
Multiple channel changes may combine to have another sign pattern.
This theorem alone does not determine stability or nonlinear safety.
Section~\ref{sec:posterclosure} checks path independence numerically at 60 frozen IEEE--39 harmonic cases; no frequency interval is certified.

Differentiating the finite formula gives
\begin{equation}
 Z_{\tau_i}(s)=s e^{-s\tau_i}u_iv_i^* . \label{eq:derivative}
\end{equation}
Away from the imaginary axis $v_i^*$ denotes the displayed analytic
row, not an additional conjugation of an analytic function.

\section{Graph meaning of the physical directions}
\status{EXACT\_IDENTITY}
Use a verified symmetric physical Laplacian only to define a basis:
$M^{-1/2}L_PM^{-1/2}=U\Lambda U^T$. The lossy model in $Z$ stays intact.
Under physical mass normalization,
\[
 \widehat{\delta D}=U^TM^{-1/2}\delta D_HM^{-1/2}U
 =\tfrac12(\widehat a\,\widehat v^*+\widehat v\,\widehat a^*),
 \quad \widehat a=U^TM^{-1/2}a,\quad \widehat v=U^TM^{-1/2}v .
\]
Thus $\widehat{\delta D}_{k\ell}
=\tfrac12(\widehat a_k\overline{\widehat v_\ell}
+\widehat v_k\overline{\widehat a_\ell})$.
Sylvester's law preserves the positive/negative counts under this
invertible congruence. Both network paths determine the modal effect;
the spatial Fourier spectrum of $\tau$ alone cannot supply its sign.
Low graph modes are not automatically the right control subspace.
The rank-two claim concerns one channel's increment, not the full
damping matrix.

\section{A graph law for inertia and aggregate frequency}
\subsection{An explicit reduced class}
Assume, separately from the full-model theorem,
\begin{equation}
 \nu=\dot q,\qquad Z_r(s)\nu=d,\qquad
 Z_r(s)=sM+\mathcal D(s)+L/s,\quad
 \mathcal D(s)=D_0+sD_1+O(s^2). \label{eq:reduced}
\end{equation}
$M$ is positive diagonal, $L=BWB^T$ is a connected undirected
Laplacian with $W>0$, and $\mathcal D$ is real analytic near zero.
For physical per-unit speed, $q=\delta/\omega_0$; consequently
$L=\omega_0 L_{\rm angle}$. $M,D_0,L$ then have units MW\,s, MW,
and MW/s. This scaling is essential.
The class (\ref{eq:reduced}) has not been derived for the full lossy
IEEE--39 PLL/DC model; singular hidden states at zero cannot be ignored.

Define
\[
 m=\one^TM\one,\quad d_0=\one^TD_0\one>0,\quad
 u_0=M^{1/2}\one/\sqrt m,\quad
 M^{-1/2}LM^{-1/2}=U\diag(0,\Lambda_\perp)U^T .
\]
Write $\widetilde D=U^TM^{-1/2}\mathcal D M^{-1/2}U$ and
$Z_\perp=sI+\widetilde D_{\perp\perp}+\Lambda_\perp/s$.
For inertia-proportional input $d=(M\one/m)p$, the COI speed
$w=\one^TM\nu/m$ satisfies
\begin{equation}
 z_{\rm coi}(s)=m[s+\widetilde D_{00}
 -\widetilde D_{0\perp}Z_\perp^{-1}\widetilde D_{\perp0}],
 \qquad w=p/z_{\rm coi}. \label{eq:coi}
\end{equation}
This is an exact Schur identity for this class wherever defined.

\begin{theorem}[Graph correction to the aggregate response moment]
Let $b_R=D_0\one-(d_0/m)M\one$ and
$b_L=D_0^T\one-(d_0/m)M\one$. Then
\begin{equation}
 z_{\rm coi}(s)=d_0+s m_{\rm app}+O(s^2),\qquad
 \boxed{m_{\rm app}=m+\one^TD_1\one-b_L^TL^\dagger b_R.} \label{eq:moment}
\end{equation}
If $D_0$ is symmetric, $b_L=b_R=b$ and
\begin{equation}
 \boxed{b^TL^\dagger b
 =m\sum_{k>0}\frac{\widetilde D_{k0}(0)^2}{\lambda_k}
 =\min_{Bf=b} f^TW^{-1}f\geq0.} \label{eq:flow}
\end{equation}
\end{theorem}
\status{PROVED\_REDUCED\_MODEL}
\begin{proof}
$Z_\perp^{-1}=s(\Lambda_\perp+s\widetilde D_{\perp\perp}+s^2I)^{-1}
=s\Lambda_\perp^{-1}+O(s^2)$. Substitute in (\ref{eq:coi}).
Both $b_L,b_R$ sum to zero. The mass-normalized spectral solution of
$Lx=b_R$ differs from $L^\dagger b_R$ only by a constant vector,
which $b_L^T$ annihilates. This identifies the cross coefficient as
$b_L^TL^\dagger b_R$. Symmetry gives the sum of squares.
For the flow identity, $f_*=WB^TL^\dagger b$ is feasible and
every feasible $f=f_*+c$ has $Bc=0$.
Its cost is $f_*^TW^{-1}f_*+c^TW^{-1}c$, proving the minimum.
\end{proof}
$m_{\rm app}$ is a response moment, not stored kinetic inertia.
It may be negative with positive physical inertia in a stable system.
The mismatch uses row/column sums of collective damping, not its
diagonal alone. Weak graph modes amplify the effect by $1/\lambda_k$.
Cycle flows belong to $\ker B$ and add the nonnegative cost above;
no separate cycle damping object is needed.
At fixed $b$, $\partial(b^TL^\dagger b)/\partial w_e
=-(B_e^TL^\dagger b)^2\leq0$.
This concerns the moment at a fixed operating point, not monotonic
improvement of nadir or stability under actual line changes.

\begin{corollary}[Negative moment forces an overshoot]
Suppose the frequency response is stable with integrable transient.
For $p(t)=-P1_{t\geq0}$, $P>0$, define signed frequency drop $x=-w$.
Then $x_\infty=P/d_0$ and
\begin{equation}
 \boxed{\int_0^\infty(x_\infty-x(t))\,dt=P m_{\rm app}/d_0^2.} \label{eq:area}
\end{equation}
If $m_{\rm app}<0$, the nadir is below the final frequency.
\end{corollary}
\begin{proof}
Expand $X(s)=P/[s z_{\rm coi}(s)]$.
The limit of $P/(d_0s)-X(s)$ at zero is the stated area.
A negative integral cannot have an everywhere nonnegative integrand.
\end{proof}
A positive moment does not rule out overshoot. A negative one does
not quantify the peak or establish violation of a declared limit.

\paragraph{Exact two-node example; no time simulation.}
Let $M=I_2$, $L=\left[\begin{smallmatrix}1&-1\\-1&1\end{smallmatrix}\right]$
in dimensionless units. Both $D_0=5I_2$ and $D_0=\diag(1,9)$ have
$m=2,d_0=10$. The former has $w/p=1/(2s+10)$, monotone.
For the latter, $b=(-4,4)^T$, $b^TL^\dagger b=16$, $m_{\rm app}=-14$,
and
\[
 \frac{w}{p}=\frac{s^2+5s+2}{2(s^3+10s^2+11s+10)}.
\]
The cubic is Hurwitz by $10,11,10>0$ and $10\cdot11>10$.
Equation (\ref{eq:area}) proves overshoot. Equal total inertia and
total damping do not determine the frequency transient.
This example is not an IEEE--39 result.

For a general deficit input $d(t)=-v1_{t\geq0}$ with $\one^Tv=P>0$,
put $a=v-(P/m)M\one$. Eliminating the disagreement input yields
\begin{equation}
 \int_0^\infty(x_\infty-x(t))\,dt
 =\frac{P m_{\rm app}}{d_0^2}+\frac{b_L^TL^\dagger a}{d_0}. \label{eq:location}
\end{equation}
The extra term comes from the numerator $P-s b_L^TL^\dagger a+O(s^2)$.
Thus the equations accept arbitrary input location. Algebraic load
events need their real input map, not an assumed mechanical step.

\section{Nonlinear causality, RoCoF and replacement}
\begin{proposition}[A causal limitation on PLL retuning]
Assume the nonlinear DDAE is well posed, its pre-event detector history
is zero, its trim is gain-independent, and gains occur only as in
(\ref{eq:pll}). At fixed other parameters, changing site $i$'s PLL
gains cannot change the physical trajectory for $0\leq t<\tau_i$.
If all tunable channels have positive delay, the trajectory before
$\tau_{\min}$ is independent of all their gains.
\end{proposition}
\status{EXACT\_IDENTITY}
\begin{proof}
The corresponding delayed error is zero on the stated interval.
Both designs have the same vector field and history there;
uniqueness of the retarded initial-value problem gives equality.
\end{proof}
Other converter states, currents, DC dynamics, SG governors and
algebraic voltages can respond immediately. This is not a claim
that GFL power is frozen before the PLL delay.
Any output-guard failure in that common interval is unrepairable
by these gains. If $\mathcal S_{\rm blind}$ is the set of $\rho$
passing all declared event/output guards on that interval, then
\[
 \{\rho:\exists K\ \text{fully feasible}\}\subseteq\mathcal S_{\rm blind},
 \qquad
 P_{GFL}^*\leq\sup_{\rho\in\mathcal S_{\rm blind}}\sum_iP_i^0\rho_i .
\]
This is an unevaluated necessary-condition relaxation, potentially
weak for short delays and long measurement windows.

For a mechanical step in (\ref{eq:reduced}), with causal controller
realization and no instantaneous synthetic-inertia feedthrough,
\[
 \dot\nu(0^+)=M^{-1}d(0^+),\qquad
 \dot w(0^+)=\one^Td(0^+)/m .
\]
An instantaneous COI rotor RoCoF limit $R_{\max}$ Hz/s therefore needs
\begin{equation}
 \sum_i2H_iS_i^0\epsilon_i
 \geq f_0|\one^Td(0^+)|/R_{\max}. \label{eq:floor}
\end{equation}
This classical condition is not novel. Minimizing $P^T\epsilon$
under it and $0\leq\epsilon\leq1$ is fractional knapsack: sort
$2H_iS_i^0/P_i^0$ decreasingly for positive dispatch and retain in
that order, with at most one final fraction. This gives only the
necessary-condition bound for that explicit input/output contract.

The frozen repository instead monitors bus-voltage phase over
$T_w=0.5$ s:
\[
 F_j(t)=\frac{\phi_j(t)-\phi_j(t-T_w)}{2\pi T_w},\quad
 R_j(t)=\frac{\phi_j(t)-2\phi_j(t-T_w)+\phi_j(t-2T_w)}{2\pi T_w^2}.
\]
Z-load events can jump algebraic voltage phase.
These metrics cannot be replaced by instantaneous rotor RoCoF.
No numerical upper replacement bound for the frozen problem is
computed in this document.

\section{The exact connection to stability and local design}
For a simple root, with nonzero eigenvalue denominator,
\[
 \lambda_p=-\frac{w^*\Delta_pv}{w^*\Delta_s v},\qquad
 \Delta_s=E+\sum_i\tau_i b_ic_i^T e^{-s\tau_i}.
\]
Delayed memory is not stored inertia. At a regular physical-port root,
\begin{equation}
 \boxed{\lambda_{\tau_i}
 =-\frac{\lambda e^{-\lambda\tau_i}
             (\ell^*u_i)(v_i^*r)}{\ell^*Z_s r}.} \label{eq:pole}
\end{equation}
This follows by using (\ref{eq:derivative}) in
$\lambda_p=-\ell^*Z_pr/(\ell^*Z_sr)$.
Phase and left/right mode overlaps matter; the harmonic sign theorem
does not replace root or nonlinear-security analysis.
The affine PLL injection also gives the consistency identity
$\lambda_{\tau_i}=-\lambda(\lambda_{\log K_{p,i}}+
\lambda_{\log K_{i,i}})$ at fixed $\rho$ and trim.
Multiple roots require subspace/nonsmooth analysis.

For $p=(\rho,\log K_p,\log K_i)$ a legitimate smooth local predictor is
\[
 \max_{\delta p}\sum_iP_i^0\delta\rho_i,\qquad
 g_j+\nabla_p g_j^T\delta p\leq0,
 \quad p+\delta p\text{ inside gain bounds and a trust region}.
\]
Spectral derivatives come from (\ref{eq:pole}); nonlinear peak
derivatives need the retarded variational equations and all active
times. This is standard sequential constrained optimization, not a
new global solver. No new optimized $\rho,K_p,K_i$ are produced here.

Nor should $m_{\rm app}\geq0$ be added as a necessary safety
constraint: overshoot below final frequency may still respect the
allowed frequency excursion. Positive harmonic damping at sampled
frequencies also does not prove stability.

\section{Exact retuning and its collective consequences}\label{sec:retuning}

\paragraph{Scope.}
The design question is whether additional replacement can be compensated
by bounded PLL gains, and which collective modes prevent that compensation.
We keep the actual nonlinear DDAE, its exact delayed linearization and
the physical ports. No new damping operator is postulated.
All delays remain externally specified. A change of delay below compares
two prescribed physical scenarios; it is not an optimization over latency.
Every formula uses a fixed state-support chart.

\subsection{The actual single-PLL return relation}
Remove the three PLL states at site $i$ and regard its angle $\theta_i$
as an input to the remaining grid. Keep all other PLLs, delays, SG,
voltage, current and DC equations. At a regular frequency write
\[
 \mathcal A_i(s;\rho)z=\mathcal B_i(s;\rho)\theta_i,\qquad
 e_i=\mathcal C_i(s;\rho)z+\mathcal D_i(s;\rho)\theta_i .
\]
The exact return relation is
\begin{equation}
 e_i=G_i(s;\rho)\theta_i,\qquad
 G_i=\mathcal D_i+\mathcal C_i\mathcal A_i^{-1}\mathcal B_i .
 \label{eq:return}
\end{equation}
The detector includes the negative local-angle term. $G_i$ is not
terminal voltage phase alone or a lossless scalar approximation.
It is independent of this site's removed $K_{p,i},K_{i,i},\tau_i$;
other controllers and delays are held fixed. Replacement dependence
includes the actual algebraic closure and equilibrium dependence.
Eliminating this PLL gives
\begin{equation}
 F_i(s)=s^2(1+t_{f,i}s)
 -(sK_{p,i}+K_{i,i})e^{-s\tau_i}G_i(s;\rho)=0 . \label{eq:scalarroot}
\end{equation}
Where $\mathcal A_i$ is nonsingular, the full determinant is
$\det(\mathcal A_i)F_i/t_{f,i}$ up to a nonzero ordering constant.
Hidden-block singularities must instead be treated with the original
characteristic operator. Not every pole need be visible in this relation.

\begin{theorem}[Exact single-site pole assignment with gain limits]
Fix $\rho$, all other controllers and delays, and
$\lambda=\alpha+\ii\omega$, $\omega>0$.
Assume $\mathcal A_i(\lambda)$ is invertible and
$G_i(\lambda;\rho)\ne0$. Define
\[
 W_i(\lambda,\rho,\tau_i)=
 \frac{\lambda^2(1+t_{f,i}\lambda)e^{\lambda\tau_i}}
 {G_i(\lambda;\rho)} .
\]
The unique real gains assigning $\lambda$ through this channel are
\begin{equation}
 \boxed{K_{p,i}=\frac{\Im W_i}{\omega},\qquad
 K_{i,i}=\Re W_i-\frac{\alpha}{\omega}\Im W_i .} \label{eq:gainexact}
\end{equation}
Under declared gain bounds, such an assignment is possible if and only if
\begin{equation}
 W_i\in\{\lambda k_p+k_i:\ k_p\in[K_p^-,K_p^+],\
                              k_i\in[K_i^-,K_i^+]\}. \label{eq:parallelogram}
\end{equation}
\end{theorem}
\status{EXACT\_IDENTITY}
\begin{proof}
Equation (\ref{eq:scalarroot}) reduces to
$\lambda K_{p,i}+K_{i,i}=W_i$.
Its imaginary and real parts give (\ref{eq:gainexact}), uniquely since
$\omega\ne0$. The image of the gain rectangle is precisely the
parallelogram in (\ref{eq:parallelogram}).
\end{proof}
This is a one-pole, one-site feasibility test with other gains frozen.
Failure is not a certificate against all distributed stabilization.
Pole assignment is established control theory; the system-specific
content is the exact return relation and the cross-mode result below.

\subsection{Transport between prescribed heterogeneous delays}
For frozen $\rho$ and an existing root $\lambda$, compare
$\tau_i$ with $\tau_i+h$. The exact gain map is
\begin{align}
 K_{p,i}(h)&=e^{\alpha h}\left[K_{p,i}\cos\omega h+
 \frac{K_{i,i}+\alpha K_{p,i}}{\omega}\sin\omega h\right],\\
 K_{i,i}(h)&=e^{\alpha h}\left[K_{i,i}\cos\omega h-
 \frac{|\lambda|^2K_{p,i}+\alpha K_{i,i}}{\omega}\sin\omega h\right].
 \label{eq:transport}
\end{align}
Indeed,
\[
 [\lambda K_{p,i}(h)+K_{i,i}(h)]e^{-\lambda(\tau_i+h)}
 =(\lambda K_{p,i}+K_{i,i})e^{-\lambda\tau_i}.
\]
The eliminated PLL equation at $\lambda$ is unchanged.
Applying this map to all prescribed heterogeneous delay changes
preserves the eliminated full-grid equations at
$\lambda,\overline\lambda$, for fixed $\rho$.
It does not ensure that these poles remain dominant or that the
other poles remain stable. All gain bounds must still hold.

The same map solves
\begin{equation}
 \frac{d}{dh}\begin{bmatrix}K_{p,i}\\K_{i,i}\end{bmatrix}
 =\begin{bmatrix}2\alpha&1\\-|\lambda|^2&0\end{bmatrix}
             \begin{bmatrix}K_{p,i}\\K_{i,i}\end{bmatrix}. \label{eq:gainflow}
\end{equation}
At $\alpha=0$, initially positive gains first lose positivity through
$K_i=0$ at $h=\omega^{-1}\arctan[K_i/(\omega K_p)]$ for increasing $h$.
Actual bounds may be reached earlier. This is a boundary of the
pole-preserving gain branch, not the maximum stable grid delay.

\begin{theorem}[Modal spillover under single-site pole preservation]
Let $\lambda=\alpha+\ii\omega$ and $\mu$ be simple characteristic
roots at the baseline, with return reductions regular at both.
Retune site $i$ by (\ref{eq:transport}) as its prescribed latency
changes, holding other sites fixed. Then
\begin{equation}
 \boxed{\left.\frac{d\mu}{dh}\right|_{0}
 =-K_{p,i}(\mu-\lambda)(\mu-\overline\lambda)
                         \frac{\partial\mu}{\partial K_{i,i}} .}
 \label{eq:spillover}
\end{equation}
The derivative on the right is in ordinary gain, not log gain.
\end{theorem}
\status{EXACT\_IDENTITY}
\begin{proof}
Differentiating (\ref{eq:scalarroot}) at a simple root gives
\[
 \mu_{K_{p,i}}=\mu\,\mu_{K_{i,i}},\qquad
 \mu_{\tau_i}=-\mu(K_{p,i}\mu_{K_{p,i}}+K_{i,i}\mu_{K_{i,i}}).
\]
Combining these with (\ref{eq:gainflow}),
\begin{align*}
 \dot\mu&=\mu_{\tau_i}
 +(K_{i,i}+2\alpha K_{p,i})\mu_{K_{p,i}}
 -|\lambda|^2K_{p,i}\mu_{K_{i,i}}\\
 &=-K_{p,i}(\mu^2-2\alpha\mu+|\lambda|^2)\mu_{K_{i,i}} .
\end{align*}
\end{proof}
Equivalently, for
$\mathcal K_i(s,h)=[sK_{p,i}(h)+K_{i,i}(h)]e^{-s(\tau_i+h)}$,
\[
 \partial_h\mathcal K_i(s,0)=
 -K_{p,i}e^{-s\tau_i}(s-\lambda)(s-\overline\lambda).
\]
For a finite change, $\mathcal K_i(s,h)-\mathcal K_i(s,0)$ has
zeros at $\lambda,\overline\lambda$, and is divisible by that
quadratic as an entire function of $s$.

\begin{corollary}[A precise single-site limitation]
If $\lambda_{K_{i,i}}\ne0$, any differentiable real retuning at
this site preserving the full complex root $\lambda$ while
$d\tau_i/dh=1$ has derivative (\ref{eq:gainflow}).
For $K_{p,i}>0$, it cannot also preserve a distinct root
$\mu\notin\{\lambda,\overline\lambda\}$ to first order when
$\mu_{K_{i,i}}\ne0$.
\end{corollary}
\begin{proof}
The gain-to-root real Jacobian is nonsingular:
$\lambda_{K_i}(\lambda\,\delta K_p+\delta K_i)$ has nonzero
real determinant for $\omega>0$, $\lambda_{K_i}\ne0$.
Compensation is therefore unique; (\ref{eq:spillover}) is nonzero.
\end{proof}
This is not an impossibility theorem for preserving only the real
parts of two roots, or for distributed retuning at other sites.
A complex root moving need not have a worsened real part.
The sign of $\Re\dot\mu$ is a prediction to test, not an assumed
destabilizing outcome.

\paragraph{Connection to Beyond Nodal Damping.}
The residue in (\ref{eq:spillover}) is
\[
 \mu_{K_{i,i}}=
 -\frac{\ell_\mu^*Z_{K_{i,i}}(\mu)r_\mu}{\ell_\mu^*Z_s(\mu)r_\mu}.
\]
The local PLL fixes the spectral polynomial; the lossy network,
inertia, P/Q and other devices fix its collective residue.
Nodal damping coefficients do not contain these left/right overlaps.
GSP can express the overlaps in a graph basis, but unsigned delay
roughness cannot substitute for their complex sign.

\subsection{Replacement at fixed physical delays}
At any prescribed $\rho$ in the regular region,
(\ref{eq:gainexact}) gives an exact candidate for one target pole.
For fixed $\lambda,\tau_i,t_{f,i}$,
\[
 W_{i,\rho_j}=-W_iG_{i,\rho_j}/G_i,\quad
 K_{p,i,\rho_j}=\Im W_{i,\rho_j}/\omega,\quad
 K_{i,i,\rho_j}=\Re W_{i,\rho_j}-\alpha\Im W_{i,\rho_j}/\omega .
\]
The derivative of $G_i$ must include equilibrium and algebraic-voltage
dependence. The sharing contract leaves $t_{f,i}$ fixed.
Holding a full complex pole is one candidate strategy.
For a requirement on real part alone, use
$\lambda=-\sigma_{\rm req}+\ii\omega$ with $\omega$ free;
fixing $\omega$ would add an unrequested design constraint.

\section{From modal prediction to a bounded replacement decision}

\begin{theorem}[Gain-bounded first-order replacement in a fixed direction]
Start at a feasible design. Fix $\rho=\rho_0+t r$, $t\ge0$,
with $P^Tr>0$, and an admissible interval $0\le t\le t_{\rm box}$.
Stack the included smooth constraint branches as $g\le0$ and put
$q=-g(p_0)\ge0$, $a=g_\rho r$, $C=g_K$.
Let the gain increments belong to a nonempty compact convex set
$\mathcal U$ containing zero. In the first-order constraint model,
\[
 \exists u\in\mathcal U:\quad t a+C u\le q
\]
if and only if
\begin{equation}
 t\,\eta^Ta\le\eta^Tq+h_{\mathcal U}(-C^T\eta)
                      \quad\text{for every }\eta\ge0 . \label{eq:dual}
\end{equation}
Consequently
\begin{equation}
 \boxed{t_{\max}^{\rm lin}
 =\min\left\{t_{\rm box},
 \inf_{\eta\ge0,\ \eta^Ta>0}
 \frac{\eta^Tq+h_{\mathcal U}(-C^T\eta)}{\eta^Ta}\right\}.}
 \label{eq:raybound}
\end{equation}
An empty infimum is $+\infty$. The maximal extra replacement
in this model is $t_{\max}^{\rm lin}P^Tr$ MW.
\end{theorem}
\status{PROVED\_REDUCED\_MODEL}
\begin{proof}
Feasibility means $q-ta\in C\mathcal U+\mathbb R_+^m$,
a closed convex set. Its separating halfspaces with finite
lower support use normals $\eta\ge0$, and
$\inf_{u\in\mathcal U}\eta^TCu=-h_{\mathcal U}(-C^T\eta)$.
This proves (\ref{eq:dual}). For $\eta^Ta\le0$ the inequality
is automatic because $q\ge0$, $0\in\mathcal U$ and $t\ge0$.
The remaining inequalities give (\ref{eq:raybound}).
\end{proof}
For the physical gain box $u_j\in[l_j,h_j]$,
$h_{\mathcal U}(z)=\sum_j\max(l_jz_j,h_jz_j)$.
The existing gain bounds are not enlarged.
If $C^T\eta=0$ for a nonnegative $\eta$, that weighted combination
of constraints has no first-order tuning authority, even if gain
increments were unbounded. Its replacement demand must be met
by baseline margin or another physical intervention.

This is standard convex separation applied to the problem; it
extends the repository's previous compensability formulation.
It is not a new general optimization theorem, a maximum over
all replacement directions, or a full nonlinear optimum.
The physical contribution sought is how the collective residues
determine the limiting combination of constraints.
For spectral maxima or trajectory peaks, retain all relevant
active branches or use directional nonsmooth analysis.
A finite list does not certify unseen roots, time instants or
disturbances.

For a uniformly verified remainder bound on the entire design region,
\[
 |g_j(p_0+(tr,u))-g_j(p_0)-t a_j-(Cu)_j|\le\varepsilon_j ,
\]
replacing $q$ by $q+\varepsilon$ gives a necessary relaxation,
and hence a regional upper bound on $t$.
Replacing $q$ by $q-\varepsilon$ in (\ref{eq:dual}) gives a
sufficient test for the included constraints. The simplified
ratio formula need not apply to $q-\varepsilon$, which can have
negative entries. No rigorous remainder bound has been obtained here.

\subsection{Analytic examples that expose the limiting mechanism}
Let $t\ge0$ denote normalized replacement and let $u\in[-\kappa,\kappa]$
be one normalized tuning direction. If the two modal constraints are
\[
 t+u\le q,\qquad t-u\le q,\qquad q\ge0,
\]
then adding them proves $t\le q$, attained with $u=0$.
Each constraint separately would allow $t=q+\kappa$.
The dual direction $\eta=(1,1)$ cancels the controller action and
detects the conflict. More tuning range does not remove this
particular first-order multimode limitation.
In contrast, if both rows are $t+u\le q$, the optimum is
$t=q+\kappa$, attained with $u=-\kappa$.
These exact algebraic examples explain what the test detects;
they do not assert that either row pattern occurs on IEEE--39.

\subsection{A sufficient gate for the complete DDE spectrum}
This standard homotopy argument prevents pole assignment from
silently ignoring the other roots. At fixed $\rho$, in a common
gauge-deflated $E=I$ state chart, let
\[
 \Delta_1(s)=\Delta_0(s)+U(s)C,\qquad
 U_i(s)=b_i^0e^{-s\tau_i^0}-b_i^1e^{-s\tau_i^1}.
\]
Both delay vectors are prescribed.
Let $\Gamma$ enclose $\Re s\ge-\sigma$, truncated by $|s|\le R$,
with $\Delta_0$ nonsingular on $\Gamma$. A sufficient condition is
\begin{equation}
 \sup_{s\in\Gamma}\|C\Delta_0(s)^{-1}U(s)\|_2<1 . \label{eq:contour}
\end{equation}
Indeed, for $0\le t\le1$ the determinant lemma and Neumann
invertibility make $\Delta_0+tUC$ nonsingular on $\Gamma$.
Its argument-principle root count is unchanged.
There are no omitted roots for $|s|>R,\Re s\ge-\sigma$ if
\[
 R>\|A_0\|_2+
 \max_{k=0,1}\sum_i\|b_i^k\|_2\|c_i\|_2e^{\sigma\tau_i^k}.
\]
Dominance of $sI$ proves this bound throughout the affine homotopy.
A zero initial root count and verified (\ref{eq:contour}) therefore
imply the full spectral margin.

The gauge must actually be removed; pinned roots on the contour
need separate deflation and are excluded here.
This sufficient condition is potentially conservative and has not
been evaluated. Sampling cannot certify its supremum.
Changing $\rho$ requires the full characteristic difference,
including $A_0,C$ and equilibrium dependence, rather than the
fixed-$\rho$ low-rank update above. Nonlinear event guards remain
separate from this spectral argument.

\paragraph{Deferred validation; no experiments in this extension.}
First check the gain-sensitivity ratio and spillover formula on
withheld simple modes. Then check finite gain transport, all
gain bounds and the complete spectrum. Only afterwards test
the frozen nonlinear events and whether a replacement predictor
improves a feasible design under identical limits. No trajectories,
optimization runs or spectral sweeps were executed here.

\section{All-PLL construction and its exact scope}\label{sec:closure}

\paragraph{Research contract.}
The primary question is whether the collective grid return can predict
an admissible PLL retuning as SG replacement increases. The proposed
mechanism is incompatibility between the gains demanded by different
oscillatory patterns. The falsifying test is that this incompatibility
fails to predict useful full-model designs, or disappears when unnecessary
pattern restrictions are removed. The consequence, if validated, is a
controller decision based on network interactions. Neighbor feedback is
a separately identified architecture extension, with its own latency and
gain contract. None of these statements assumes a strict advantage exists.

\subsection{Open every PLL: the plant remains the full grid}
Remove all PLL triples $(\theta,\omega_p,\xi)$ while keeping the $n$
angles as external inputs to the other device/network equations:
\[
 \mathcal A(s;\rho)z=\mathcal B(s;\rho)\theta,\qquad
 e=\mathcal C(s;\rho)z+\mathcal D(s;\rho)\theta .
\]
For $\mathcal A$ nonsingular, define
\begin{equation}
 G(s;\rho)=\mathcal D+\mathcal C\mathcal A^{-1}\mathcal B,
 \qquad D_f(s)=s^2(I+sT_f). \label{eq:allreturn}
\end{equation}
Unlike $G_i$ in (\ref{eq:return}), $G$ is independent of \emph{all}
removed PLL gains and detector delays. The current, DC, SG, lossy AC,
P/Q and equilibrium dependencies remain in $G$. This independence uses
the actual model contract: gains enter only the removed PLL equations.
Do not silently reuse it for a model with other gain-dependent channels.

For local diagonal gains $K_p,K_i$, the exact reduced characteristic is
\begin{equation}
 F(s;\rho,K)=D_f(s)-(sK_p+K_i)E_\tau(s)G(s;\rho),
 \qquad E_\tau=\diag(e^{-s\tau_i}). \label{eq:allF}
\end{equation}
Its determinant describes full-model roots wherever the eliminated
block is regular; hidden-block singularities need the original descriptor.
No approximation of the exponentials or modal diagonalization is used.

\begin{theorem}[Analytical gains for a prescribed collective pattern]
Fix a replacement $\rho$, $\lambda=\alpha+\ii\omega$ with $\omega>0$,
and a nonzero PLL-angle vector $q\in\mathbb C^n$. Assume
$\mathcal A(\lambda;\rho)$ is nonsingular and
$y_i=[G(\lambda;\rho)q]_i\ne0$ for every site.
Set
\begin{equation}
 W_i=\frac{\lambda^2(1+t_{f,i}\lambda)e^{\lambda\tau_i}q_i}{y_i},
 \qquad
 \boxed{K_{p,i}=\frac{\Im W_i}{\omega},\quad
 K_{i,i}=\Re W_i-\alpha\frac{\Im W_i}{\omega}.} \label{eq:allgain}
\end{equation}
These are the unique real local gains for which $F(\lambda)q=0$.
They are admissible precisely when they obey the declared gain bounds.
\end{theorem}
\status{EXACT\_IDENTITY}
\begin{proof}
Row $i$ of $F(\lambda)q=0$ is
$(\lambda K_{p,i}+K_{i,i})e^{-\lambda\tau_i}y_i
=\lambda^2(1+t_{f,i}\lambda)q_i$.
Its two real equations give (\ref{eq:allgain}). Reconstruct
$z=\mathcal A^{-1}\mathcal Bq$, $\omega_p=\lambda q$ and
$\xi=K_iE_\tau y/\lambda$ to recover the full eigenvector.
Real model coefficients also give the conjugate eigenpair.
\end{proof}
If $y_i=0$ and the right side is nonzero, assignment is impossible.
If both vanish, that row imposes no gain equation and must be treated
without division. Nonzero $q$ is required; its complex scale cancels
in (\ref{eq:allgain}). This result assigns an eigenpair, not dominance,
simple multiplicity, the complete spectrum, or nonlinear safety.
Conversely, any existing local-gain design with such an eigenpair
and nonzero responses satisfies this same formula. Thus allowing
the pole and pattern to vary covers this regular eigenpair chart;
holding a convenient pattern fixed is the additional restriction.

The collective dependence is explicit:
\[
 y_i=\sum_jG_{ij}(\lambda;\rho)q_j .
\]
Thus even local gains depend on the spatial amplitudes and phases of
the whole pattern. Replacing this sum by $G_{ii}q_i$ is an additional
model approximation. Offline use of $G$ does not require online
communication or make the controller distributed.

\subsection{Exact gain compatibility for several patterns}
For $m$ distinct nonreal target pairs $(\lambda_r,q_r)$, construct
$W_{ir}$ from (\ref{eq:allgain}). Let
\[
 C_\Lambda=
 \begin{bmatrix}
 \alpha_1&1\\ \omega_1&0\\
 \vdots&\vdots\\ \alpha_m&1\\ \omega_m&0
 \end{bmatrix},\qquad
 w_i=\begin{bmatrix}\Re W_{i1}\\\Im W_{i1}\\
 \vdots\\\Re W_{im}\\\Im W_{im}\end{bmatrix}.
\]
Then, with no gain bounds, simultaneous local assignment is possible
if and only if
\begin{equation}
 \boxed{(I-C_\Lambda C_\Lambda^\dagger)w_i=0
 \quad\hbox{for all }i.} \label{eq:compat}
\end{equation}
The gains are $[K_{p,i},K_{i,i}]^T=C_\Lambda^\dagger w_i$.
Since $\omega_r>0$, $C_\Lambda$ has rank two and the gains are unique;
bounded feasibility also requires membership in the gain rectangle.
\status{EXACT\_IDENTITY}
The proof is the range criterion for $C_\Lambda k_i=w_i$.
For two pairs this is four real equations for two real gains at each
site. The conflict is therefore testable, rather than inferred from
large or negative damping coefficients.

\textbf{Scope restriction.} Failure proves infeasibility for these
\emph{specified poles and patterns}. It does not prove that local
controllers cannot stabilize the same replacement with other modes.
It is not a decentralized fixed-mode theorem. Patterns and target
frequencies must be allowed to vary before claiming an architectural
stabilization limit.

\subsection{Which additional measurement can remove a conflict?}
For an admissible directed information pattern $\mathcal N_i$, including
the local signal, consider the explicit extension
\begin{align*}
 t_{f,i}\dot\omega_{p,i}
 &=\xi_i-\omega_{p,i}
   +\sum_{j\in\mathcal N_i}k_{p,ij}e_j(t-\tau_j-d_{ij}),\\
 \dot\xi_i&=\sum_{j\in\mathcal N_i}k_{i,ij}e_j(t-\tau_j-d_{ij}),
 \qquad \dot\theta_i=\omega_{p,i}.
\end{align*}
Every $d_{ij}\ge0$ is an exogenous communication delay, with $d_{ii}=0$.
Signals are calibrated detector errors with declared units, sign and
timestamps. Their zero locked values preserve the equilibrium.
Off-diagonal gains and their bounds define a new controller contract;
they must not be treated as the old twenty-gain architecture.

For fixed target $(\lambda_r,q_r)$ define
\[
 y_r=G(\lambda_r;\rho)q_r,\quad
 z_{ijr}=e^{-\lambda_r(\tau_j+d_{ij})}(y_r)_j,\quad
 b_{ir}=\lambda_r^2(1+t_{f,i}\lambda_r)(q_r)_i .
\]
The row condition is
\[
 \sum_{j\in\mathcal N_i}
  [\lambda_r z_{ijr}k_{p,ij}+z_{ijr}k_{i,ij}]=b_{ir}.
\]
Stack real and imaginary parts into the real linear system
\begin{equation}
 A_i x_i+B_i u_i=b_i. \label{eq:neighbor}
\end{equation}
$A_i$ contains the two self-gain columns; $B_i$ contains the proposed
neighbor-gain columns. A proportional column for $j$ is the stacking
of $\Re(\lambda_r z_{ijr}),\Im(\lambda_r z_{ijr})$; an integral column
stacks $\Re z_{ijr},\Im z_{ijr}$. Apply a fixed declared row scaling
before using residual norms or singular values. Exact feasibility is
unchanged by invertible row scaling.

\begin{theorem}[Compatibility and minimum added gain]
For (\ref{eq:neighbor}), let
$\Pi_i=I-A_iA_i^\dagger$, $r_i=\Pi_i b_i$ and $H_i=\Pi_i B_i$.
Assume initially that self gains are unrestricted.
There exist real gains realizing every prescribed pair at site $i$
if and only if $r_i\in\operatorname{range}H_i$.
With added-gain metric $R_i\succ0$, the unique minimum-metric neighbor
gain is
\begin{align}
 u_i^\star&=R_i^{-1}H_i^T
   (H_iR_i^{-1}H_i^T)^\dagger r_i, \label{eq:neighborgain}\\
 \kappa_{i,\min}^2&=
   r_i^T(H_iR_i^{-1}H_i^T)^\dagger r_i. \label{eq:neighborcost}
\end{align}
Self gains then solve $A_ix_i=b_i-B_iu_i^\star$.
\end{theorem}
\status{EXACT\_IDENTITY}
\begin{proof}
Project (\ref{eq:neighbor}) onto the orthogonal complement of
$\operatorname{range}A_i$. This gives $H_iu_i=r_i$, which is also
sufficient because the remaining residual belongs to
$\operatorname{range}A_i$. Whiten $u_i$ with $R_i^{1/2}$ and apply
the minimum-norm solution of a consistent linear system; this proves
both formulas. Its nullspace is orthogonal in the metric to the
displayed solution, proving minimality and uniqueness.
\end{proof}
An added budget $u_i^TR_iu_i\le\kappa_i^2$ is feasible if and only if
the range test holds and $\kappa_{i,\min}\le\kappa_i$, still with
unrestricted self gains. With bounds on either gain family, solve
the exact convex feasibility/QP problem (\ref{eq:neighbor}) with
those bounds; the displayed unconstrained minimizer is not sufficient.
Polynomial graph coefficients shared across sites require stacking
all rows into one linear problem, rather than independent row solves.

With $d_i=|\mathcal N_i|$ incoming signals, the assignment matrix has
$2d_i$ columns and $2m$ rows. Arbitrary modal demands require
$d_i\ge m$ and full row rank. One neighbor can resolve a two-pair
conflict only when its two columns supply the missing directions;
communication delays can destroy or improve this rank/conditioning.
A graph edge count alone is insufficient. These statements concern
prescribed eigenstructure, not global stabilizability.

\subsection{Why arbitrary powers of a Laplacian do not close the problem}
Let $L$ be a normalized, symmetric, connected Laplacian and $K(L)$
a polynomial graph filter. For delayed source measurements,
\begin{equation}
 [L,K(L)E_\tau(s)]=K(L)[L,E_\tau(s)]. \label{eq:polycomm}
\end{equation}
If $K(L)$ is invertible, the left side vanishes exactly when the
rightmost commutator vanishes. Its entry is
$L_{ij}(e^{-s\tau_j}-e^{-s\tau_i})$.
Commutation throughout a frequency interval therefore forces equal
delays along all edges, hence uniform delay on a connected graph.
At isolated frequencies, phase coincidence can be an exception.
Furthermore,
\begin{equation}
 \|[L,E_\tau(\ii\Omega)]\|_F^2
 =8\sum_{\{i,j\}\in\mathcal E}w_{ij}^2
       \sin^2\!\left(\frac{\Omega(\tau_i-\tau_j)}2\right).
 \label{eq:freqcomm}
\end{equation}
\status{EXACT\_IDENTITY}
Both formulas follow entrywise; each undirected edge occurs twice.
If $\sigma_{\min}(K(L))\ge\gamma>0$, the remaining commutator norm
is at least $\gamma\|[L,E_\tau]\|_F$.
This is a channel-mixing result, not a stability lower bound.
It does not apply unchanged to nonuniform diagonal gains plus
graph terms, or to general edge-delayed filters.

The physical network basis $U$ is useful for expressing $q=Ua$ and
decomposing $Gq$. Using all graph coordinates is an exact basis
change; truncating $a$ restricts the design and needs separate evidence.
The actual model need not commute with $L$, even at zero delay.
Powers $L^\ell$ encode weighted walks, not an independent physical
cycle basis. No new instability is established by such powers alone.

\subsection{A calculable replacement predictor and corrector}
Use (\ref{eq:allgain}) as an exact gain-elimination map
$K=\Phi(\rho,\lambda,q)$ on a regular chart. Normalize one nonzero
complex component of $q$ to one. For fixed delays and nonzero
$q_i,y_i,W_i$, its differential is
\begin{align}
 dW_i&=W_i\left[
 \left(\frac2\lambda+\frac{t_{f,i}}{1+t_{f,i}\lambda}+\tau_i\right)d\lambda
 +\frac{dq_i}{q_i}-\frac{dy_i}{y_i}\right],\\
 dy&=G\,dq+G_s q\,d\lambda+
             \sum_jG_{\rho_j}q\,d\rho_j,\\
 dK_{p,i}&=\frac{\Im(dW_i)-K_{p,i}\,d\omega}{\omega},\\
 dK_{i,i}&=\Re(dW_i)-\alpha\,dK_{p,i}-K_{p,i}\,d\alpha .
 \label{eq:gaindifferential}
\end{align}
Use the undivided quotient rule at a zero numerator. Derivatives of
$G$ include the algebraic closure and any equilibrium dependence.
The chain rule for another simple root is
\[
 d\mu=\sum_j\mu_{\rho_j}d\rho_j+
 \sum_i(\mu_{K_{p,i}}dK_{p,i}+\mu_{K_{i,i}}dK_{i,i}),
 \qquad
 \mu_p=-\frac{\ell^*Z_p r}{\ell^*Z_s r},
\]
where the physical-port representation must be regular; otherwise
use the full descriptor sensitivity. This retains the collective,
reactive and delay effects, including the spillover in
(\ref{eq:spillover}).

In normalized real coordinates
$\zeta=(\rho,\alpha,\omega,\Re q_{\rm free},\Im q_{\rm free})$, let
$p(\zeta)=(\rho,\Phi(\rho,\lambda,q))$.
The replacement problem on this chart is
\[
 \min_\zeta\ \sum_iP_i^0(1-\rho_i),\qquad
 g_j(p(\zeta))\le0,\quad \alpha\le-\sigma_{\rm req},
\]
with the actual gain/replacement bounds, every required physical root,
and all nonlinear event guards included in $g$. This is an exact
reparameterization on the declared chart, not a convexification of
the security constraints. A computable local subproblem is
\begin{equation}
 \max_{\delta\zeta}\ (P^0)^T\delta\rho
       -\frac{\beta}{2}\delta\zeta^T R_\zeta\delta\zeta,\quad
 g_k+J_{g,k}J_{p,k}\delta\zeta\le0,\quad
 \delta\zeta\in\mathcal T_k . \label{eq:closureQP}
\end{equation}
Here $\beta>0$, $R_\zeta\succ0$ and the convex trust region
$\mathcal T_k$ are declared in advance. It is a convex QP when
$\mathcal T_k$ is polyhedral. Include all active smooth branches of
peak/spectral constraints; a single gradient of their maximum is
not generally valid. $J_p$ uses (\ref{eq:gaindifferential}).
This is a standard SQP predictor with an analytical gain map.

The method to implement after the algebra gate is:
\begin{enumerate}
 \item Start from a fully feasible design and a simple visible oscillatory
 root, with its actual PLL pattern. Confirm reconstruction by
 (\ref{eq:allgain}); this parity test is mandatory.
 \item At a frozen $\rho$-direction, compute the predicted gain change
 from (\ref{eq:gaindifferential}). Alternatively allow $\rho$ to vary
 in a trust region and maximize $(P^0)^T d\rho$ using all active
 linearized security constraints after substituting $dK$.
 \item Permit $\omega$, the admissible real part, and the complex
 pattern to vary. Fixing them is a declared restricted predictor,
 not a fair upper bound on local-controller capability.
 \item Evaluate a finite trial through $K=\Phi(\rho,\lambda,q)$,
 rather than applying only the tangent gain update.
 This makes the chosen eigenpair exact within arithmetic/model error.
 Reject bound violations and singular or poorly conditioned charts.
 \item Check the complete required spectral region and the frozen
 nonlinear guards. If another mode becomes limiting, include it and
 use the compatibility equations with variable patterns in a local
 Newton/SQP correction. Reduce the step when predicted and actual
 constraint changes disagree.
 \item Stop at a declared local stationarity/residual criterion,
 a chart failure, or an unresolved feasibility conflict. Record which.
 Do not report a restricted-family stop as the full replacement maximum.
\end{enumerate}
For prescribed patterns, the gain solve is explicit or convex. Jointly
varying patterns and replacement makes the outer problem nonconvex.
With all $q$ coordinates free, this parameterization does not promise
dimension reduction or computational superiority. Local Newton
convergence requires a nonsingular normalized residual Jacobian;
root collisions require a subspace/nonsmooth method. No global
convergence or quantitative speedup is asserted.

No frequency, RoCoF, voltage or actuator guard is inferred merely from
pole assignment. Their original nonlinear definitions and all frozen
events remain mandatory at the later validation stage. A sufficient
full-spectrum continuation condition was given above; it has not been
evaluated. Exact local gains alone are not an exact security frontier.

\subsection{Evidence gates and literature boundary}
Pole assignment with delay, eigenstructure design, sparse feedback and
spectral optimization are established methods
\cite{ram,katewa,vanbiervliet,mao}.
Ram et al.\ treat delayed feedback in symmetric mechanical systems;
Katewa and Pasqualetti treat sparse static state-feedback assignment;
Vanbiervliet et al.\ treat nonsmooth spectral optimization for retarded
systems, including continuation of nonlinear equilibria.
These antecedents preclude claiming that the matrix solve, projection,
or iterative optimization is itself new. The narrower candidate is
the bounded PLL compatibility law under the full-grid return and its
verified consequence for replacement or latency decisions.
Recent converter-model comparisons reinforce the need to retain
hidden dynamics \cite{coffey}; supplementary GFL power-oscillation
controllers \cite{chakraborty} are a different action set from PLL-only tuning.

Before new grid simulations, close three gates:
(i) the exact gain map, compatibility and derivatives pass algebra checks;
(ii) the controller architecture, fixed delays, bounds, target-selection
rule and falsification criteria are written down;
(iii) the first validation has a numerical prediction to test, not a
desired replacement percentage. Gate (i) is mathematical only; IEEE--39
parity, all-root checks and event validation were the subsequent gates.
Section~\ref{sec:mapvalidation} reports their targeted execution.

The preregistered first experiment tests whether the formula reconstructs known
gains and predicts a finite replacement correction on withheld modes.
Only after that succeeds should a matched-budget design comparison run.
Neighbor feedback is investigated only if the local modal conflict is
material after relaxing arbitrary pattern choices. A rigorous claim
that neighbors enable otherwise impossible secure replacement would
need a valid local-architecture upper bound below a fully feasible
distributed design; fixed-pattern incompatibility is not that bound.


% BEGIN GENERATED ALL-PLL VALIDATION
\begin{thebibliography}{99}
\bibitem{pates} R. Pates and E. Mallada,
\emph{Damping, Inertia, and Delay Robustness Trade-offs in Power Systems},
MTNS, 2018.
\url{https://mallada.ece.jhu.edu/pubs/2018-MTNS-PM.pdf}.
\bibitem{paganini} F. Paganini and E. Mallada,
\emph{Global performance metrics for synchronization of heterogeneously
rated power systems: The role of machine models and inertia},
arXiv:1710.07195v2, 2017.
\url{https://arxiv.org/html/1710.07195v2}.
\bibitem{pagnier} L. Pagnier and P. Jacquod,
\emph{Optimal placement of inertia and primary control: a matrix
perturbation theory approach}, arXiv:1906.06922, 2019.
\url{https://arxiv.org/abs/1906.06922}.
\bibitem{huang} L. Huang, H. Xin, W. Dong, and F. D\"orfler,
\emph{Impacts of Grid Structure on PLL-Synchronization Stability of
Converter-Integrated Power Systems}, arXiv:1903.05489, 2019.
\url{https://arxiv.org/abs/1903.05489}.
\bibitem{wilson} J. Wilson, A. Nelson, and B. Farhang-Boroujeny,
\emph{Parameter Derivation of Type-2 Discrete-Time Phase-Locked Loops
Containing Feedback Delays}, IEEE TCAS-II, 56(12):886--890, 2009.
\url{https://doi.org/10.1109/TCSII.2009.2034197}.
\bibitem{zitek} P. Z\'itek, J. Fi\v{s}er, and T. Vyhl\'idal,
\emph{Dimensional analysis approach to dominant three-pole placement
in delayed PID control loops}, Journal of Process Control,
23(8):1063--1074, 2013.
\url{https://www.sciencedirect.com/science/article/abs/pii/S0959152413001236}.
\bibitem{zhang} W. Zhang, Y. Cui, and X. Ding,
\emph{An Improved Analytical Tuning Rule of a Robust PID Controller
for Integrating Systems with Time Delay Based on the Multiple
Dominant Pole-Placement Method}, Symmetry, 12(9):1449, 2020.
\url{https://doi.org/10.3390/sym12091449}.
\bibitem{ram} Y. M. Ram, J. E. Mottershead, and M. G. Tehrani,
\emph{Partial pole placement with time delay in structures using the
receptance and the system matrices}, Linear Algebra and Its Applications,
434(7):1689--1696, 2011.
\url{https://doi.org/10.1016/j.laa.2010.07.014}.
\bibitem{katewa} V. Katewa and F. Pasqualetti,
\emph{Minimum-Gain Pole Placement With Sparse Static Feedback},
IEEE Transactions on Automatic Control, 66(8):3445--3459, 2021.
\url{https://doi.org/10.1109/TAC.2020.3018615}.
\bibitem{vanbiervliet} J. Vanbiervliet, K. Verheyden, W. Michiels,
and S. Vandewalle, \emph{A nonsmooth optimisation approach for the
stabilisation of time-delay systems}, ESAIM: COCV, 14(3):478--493, 2008.
\url{https://doi.org/10.1051/cocv:2007060}.
\bibitem{mao} X. Mao and H. Dai,
\emph{Partial eigenvalue assignment with time delay robustness},
Numerical Algebra, Control and Optimization, 3(2):207--221, 2013.
\url{https://doi.org/10.3934/naco.2013.3.207}.
\bibitem{coffey} S. Coffey, S. Harrison, A. Egea-Alvarez, and L. Xu,
\emph{Assessing the Capability and Limitations of Small Signal
Simplified Grid Following Models for the Identification of
Sub-Synchronous Oscillations}, IEEE Access, 2025.
\url{https://doi.org/10.1109/ACCESS.2025.3586197}.
\bibitem{chakraborty} T. Chakraborty, A. Pal, and S. Maleki,
\emph{Damping LFOs: Grid Following with Power Oscillation Damping
vs. Grid Forming vs. PSS}, arXiv:2505.24204, 2025.
\url{https://arxiv.org/abs/2505.24204}.

\bibitem{beyn} W.-J. Beyn,
\emph{An integral method for solving nonlinear eigenvalue problems},
arXiv:1003.1580, 2010.
\url{https://arxiv.org/abs/1003.1580}.
\bibitem{lessard} J.-P. Lessard and J. D. Mireles James,
\emph{A functional analytic approach to validated numerics for
eigenvalues of delay equations}, Journal of Computational Dynamics,
7(1):123--158, 2020.
\url{https://doi.org/10.3934/jcd.2020005}.
\bibitem{michielsdesign} W. Michiels and S. Gumussoy,
\emph{Eigenvalue based algorithms and software for the design of
fixed-order stabilizing controllers for interconnected systems with
time-delays}, in Delay Systems: Advances in Delays and Dynamics,
Springer, 2014; author version arXiv:2003.05496.
\url{https://doi.org/10.1007/978-3-319-01695-5_18}.

\bibitem{markovic} U. Markovi\'c, O. Stanojev, P. Aristidou,
E. Vrettos, D. Callaway, and G. Hug,
\emph{Understanding Small-Signal Stability of Low-Inertia Systems},
IEEE Transactions on Power Systems, 36(5):3997--4017, 2021.
\url{https://doi.org/10.1109/TPWRS.2021.3061434}.
\bibitem{chuteng} Z. Chu and F. Teng,
\emph{Stability Constrained Optimization in High IBR-Penetrated
Power Systems---Part I: Constraint Development and Unification},
author version, arXiv:2307.12151.
\url{https://arxiv.org/abs/2307.12151}.
\bibitem{jia} H. Jia et al.,
\emph{Converter-Driven Stability Constrained Generation Expansion
for High IBG-Penetrated Power Systems},
IEEE Transactions on Power Systems, 2025.
\url{https://doi.org/10.1109/TPWRS.2025.3558879}.
\bibitem{peaucelle} D. Peaucelle and Y. Ebihara,
\emph{Robust stability analysis of discrete-time systems with
parametric and switching uncertainties}, IFAC World Congress, 2014.
\url{https://skoge.folk.ntnu.no/prost/proceedings/ifac2014/media/files/0282.pdf}.
\bibitem{bindel} D. Bindel and A. Hood,
\emph{Localization Theorems for Nonlinear Eigenvalue Problems},
SIAM Journal on Matrix Analysis and Applications, 2013.
\url{https://www.cs.cornell.edu/~bindel/papers/2013-simax.pdf}.
\bibitem{borgioli} F. Borgioli, W. Michiels, D. Lu, and B. Vandereycken,
\emph{A globally convergent method to compute the real stability
radius for time-delay systems}, author manuscript, 2019.
\url{https://www.unige.ch/math/vandereycken/papers/published_Borgioli_MLV.pdf}.
\bibitem{arb} F. Johansson,
\emph{Arb: Efficient Arbitrary-Precision Midpoint-Radius Interval
Arithmetic}, IEEE Transactions on Computers, 66(8):1281--1292, 2017.
\url{https://doi.org/10.1109/TC.2017.2690633}.
\end{thebibliography}
\end{document}
